% supplement_vn.tex -- the supplementary document of paper_v2_vn (Tài liệu bổ sung, N31): the
% former appendix of the paper (sections/appendix.tex, with tables t2, t3b, t5, t6, t8, t9 and
% figure f1), built as a document of its own. Its sections are lettered A-H and its tables and
% figure numbered S1, S2, ... Numbers without the S (Mục 4, Bảng 3) are the paper's: they are read
% from paper_v2_vn.aux with the xr package (labels prefixed P-), so the paper is built first. The
% latexmkrc in this folder does that on its own, also on Overleaf (Menu, Main document:
% supplement_vn.tex). The preamble is the paper's, copied: keep the two alike.
% Compile with LuaLaTeX and BibTeX, like the paper. It uses the same three .bib files.
\RequirePackage{iftex}
\ifLuaTeX\else
  \errmessage{supplement_vn.tex needs LuaLaTeX. In Overleaf: Menu, Compiler, LuaLaTeX}%
  \expandafter\stop
\fi
\PassOptionsToPackage{quiet}{fontspec}   % notes on font features go to the log, not the warnings
\PassOptionsToPackage{dvipsnames}{xcolor}
\PassOptionsToPackage{obeyspaces,spaces}{url}   % \code keeps its spaces and can break at them
% nonacm: a draft not yet at an ACM venue, so no conference line, copyright block or ACM reference
% box.
\documentclass[sigconf,nonacm]{acmart}
\makeatletter
\@ACM@balancefalse   % no column balancing: acmart's \balance would come in the second column
\makeatother         % of the last page and only warn that it cannot balance there
\raggedbottom

% Vietnamese is the main language: babel's captions (Hình, Bảng, Phụ lục) and line breaking.
% English stays loaded for the reference list, which keeps the language of the works (N7), for
% ACM's CCS terms, and for the English words of quotations (\foreignlanguage{english}, N13).
\usepackage[english,vietnamese]{babel}
\providecommand\ccsname{CCS Concepts}
\addto\captionsvietnamese{%
  \renewcommand\abstractname{Tóm tắt}%          babel's is Tóm tắt nội dung
  \renewcommand\keywordsname{Từ khóa}%
  \renewcommand\acksname{Lời cảm ơn}%
  \renewcommand\refname{Tài liệu tham khảo}%    babel's is only Tài liệu
  \renewcommand\ccsname{Phân loại CCS}}
% acmart prints the heading of the CCS concepts as fixed text; this makes it \ccsname, and sets
% ACM's terms under it as English.
\usepackage{etoolbox}
\makeatletter
\patchcmd{\maketitle}{\@specialsection{CCS Concepts}%
         \noindent\@concepts\par}{\@specialsection{\ccsname}%
         \noindent\foreignlanguage{english}{\@concepts}\par}{}%
  {\PackageWarningNoLine{supplement_vn}{acmart's CCS block has changed, so its heading stays
     English; see \string\ccsname}}   % the same code in acmart 2.03 (2024) and 2.21 (2026)
\makeatother
\setlength{\emergencystretch}{1em}   % Vietnamese words are not hyphenated: a little more stretch
\AtEndPreamble{\hypersetup{pdflang={vi}}}   % the PDF's language; runs after acmart's pdflang={en}

\usepackage{array}
\usepackage{tabularx}
\usepackage{tikz}
\usetikzlibrary{arrows.meta,positioning}
\usepackage{pgfplots}
\pgfplotsset{compat=1.18}
\pgfkeys{/pgf/number format/use comma}   % axis numbers as 0,05 and 1.000, as in numbers.tex
\usepackage{icomma}   % in a formula, a comma between digits is a decimal comma: $-0,034$

\usepackage{xurl}   % long links may break anywhere in the narrow columns
% File names and meta-paths in typewriter type; they may break only after / _ . or at a space, so
% a meta-path such as PUPUP is never split.
\DeclareUrlCommand\codeurl{\def\UrlBreaks{\do\/\do\_\do\.}\def\UrlBigBreaks{}\urlstyle{tt}}
\DeclareRobustCommand{\code}[1]{\codeurl{#1}}
\newcolumntype{L}{>{\raggedright\arraybackslash}X}
% Wide tables and figures span both columns.
\newenvironment{widetable}[1][tbp]{\begin{table*}[#1]}{\end{table*}}
\newenvironment{widefigure}[1][tbp]{\begin{figure*}[#1]}{\end{figure*}}
% Figure colours, as in paper_v2: searched sets blue, manual orange, random draws aqua.
\definecolor{cSearch}{HTML}{2A78D6}
\definecolor{cManual}{HTML}{EB6834}
\definecolor{cRandom}{HTML}{1BAF7A}
\definecolor{cOther}{HTML}{7C7B76}
% Every number of the paper: \val{key}, from numbers.tex (Vietnamese number format, N4).
% numbers.tex -- written by make_numbers.py; do not edit by hand, run the script again.
% 2026-11-10: python make_numbers.py --lmsps LMSPS --han han-improve --lang vi
% The raw-file counts come from extraction_report.json; the block below names the copy used.
% Vietnamese number format (--lang vi): 0,666 and 1.147.920; versions, commits and path sets as in
% English. The notes after the values stay in English.
% \val{key} prints a number; an unknown key prints ??key in bold and raises a LaTeX warning.
\makeatletter
\providecommand{\setval}[2]{\@namedef{val@#1}{#2}}
\DeclareRobustCommand{\val}[1]{\@ifundefined{val@#1}{\textbf{??#1}%
  \PackageWarning{numbers}{No value `#1' in numbers.tex}}{\@nameuse{val@#1}}}
\makeatother

% --- Data: the raw file, the fixes and the clean graph (LMSPS notes, Parts 2-4) ------------------
\setval{raw-records}{548.552}  % extraction_report.json (review/evidence/extraction_report.json)
\setval{raw-usable}{542.684}  % extraction_report.json (review/evidence/extraction_report.json)
\setval{raw-reviews}{7.593.244}  % extraction_report.json (review/evidence/extraction_report.json)
\setval{raw-categories}{49.732}  % extraction_report.json (review/evidence/extraction_report.json)
\setval{raw-pos-pct}{78,6}  % audit_before_fix.json: share of review lines rated 4 or 5
\setval{raw-unhelpful-pct}{1,35}  % positive reviews dropped as unhelpful (>= 5 votes, < 20% helpful)
\setval{placeholder-ids}{6}  % extraction_report.json (review/evidence/extraction_report.json)
\setval{placeholder-lines}{1.147.920}  % extraction_report.json (review/evidence/extraction_report.json)
\setval{placeholder-main-lines}{945.065}  % extraction_report.json (review/evidence/extraction_report.json): ATVPDKIKX0DER
\setval{placeholder-second-lines}{201.770}  % extraction_report.json (review/evidence/extraction_report.json): A3UN6WX5RRO2AG
\setval{edition-pairs}{64.814}  % extraction_report.json (review/evidence/extraction_report.json)
\setval{edition-clusters}{30.648}  % extraction_report.json (review/evidence/extraction_report.json)
\setval{edition-products}{72.280}  % extraction_report.json (review/evidence/extraction_report.json)
\setval{raw-works}{501.052}  % usable products after merging editions
\setval{kcore-iterations}{13}  % extraction_report.json (review/evidence/extraction_report.json)
\setval{catpath-rows}{2.288.922}  % extraction_report.json (review/evidence/extraction_report.json)
\setval{placeholder-main-max}{1.056}  % most reviews on one product
\setval{placeholder-second-max}{1.193}
\setval{placeholder-found}{4}  % IDs found by the rule, not named in advance
\setval{placeholder-found-max-lo}{57}  % their most reviews on one product
\setval{placeholder-found-max-hi}{149}
\setval{placeholder-found-lines}{1.085}
\setval{edition-size-min}{2}  % products in an edition cluster
\setval{edition-size-max}{24}
\setval{edition-size-two-pct}{79,5}  % clusters of exactly two products
\setval{rule-pos-rating}{4}  % a positive review is rated at least this
\setval{rule-helpful-votes}{5}  % a review with at least this many votes ...
\setval{rule-helpful-pct}{20}  % ... and fewer helpful ones than this is dropped
\setval{rule-kcore}{5}
\setval{rule-edition-shared}{3}  % identical review lines two editions must share ...
\setval{rule-edition-overlap-pct}{50}  % ... covering this share of the smaller one
\setval{rule-edition-maxrun}{50}  % review lines on more products than this are ignored
\setval{rule-placeholder-repeats}{50}  % reviews on one product that mark a placeholder
\setval{before-products}{90.659}
\setval{before-categories}{39.754}
\setval{before-users}{171.285}
\setval{before-edges}{2.612.766}
\setval{before-maxdeg}{55.572}  % largest user degree
\setval{before-meddeg}{7}  % median user degree
\setval{before-clusters}{17.726}  % duplicate-edition clusters
\setval{before-clustered}{44.591}  % products in those clusters
\setval{before-clustered-pct}{49,2}
\setval{before-copyedges-pct}{55,0}  % edges that copy an interaction already in the graph
\setval{before-valleak-pct}{72,9}  % validation edges with a training copy
\setval{before-testleak-pct}{73,0}  % test edges with a training copy
\setval{before-valtitle-pct}{60,3}  % validation edges whose user has a training edge to a same-title product
\setval{before-testtitle-pct}{60,2}  % the same for test edges
\setval{before-alsobought}{188.257}  % distinct also-bought product pairs
\setval{after-products}{48.883}
\setval{after-categories}{37.411}
\setval{after-users}{53.569}
\setval{after-edges}{819.847}
\setval{after-maxdeg}{2.134}  % largest user degree
\setval{after-meddeg}{8}  % median user degree
\setval{after-clusters}{6}  % duplicate-edition clusters
\setval{after-clustered}{12}  % products in those clusters
\setval{after-clustered-pct}{0,025}
\setval{after-copyedges-pct}{0,007}  % edges that copy an interaction already in the graph
\setval{after-valleak-pct}{0,011}  % validation edges with a training copy
\setval{after-testleak-pct}{0,010}  % test edges with a training copy
\setval{after-valtitle-pct}{0,86}  % validation edges whose user has a training edge to a same-title product
\setval{after-testtitle-pct}{0,80}  % the same for test edges
\setval{after-alsobought}{87.901}  % distinct also-bought product pairs
\setval{after-zerocat-pct}{60,7}  % categories with an all-zero name vector
\setval{after-leafcats-med}{5}  % leaf categories per product
\setval{title-coverage-pct}{95,8}  % products with a non-zero title vector
\setval{title-dim}{256}
\setval{title-vocab}{12.699}
\setval{cat-vocab}{128}
\setval{user-bits}{8}
\setval{user-dim}{64}
\setval{rel-pp-rows}{121.075}  % also-bought rows per direction
\setval{rel-cc}{37.400}
\setval{rel-pc}{427.565}
\setval{rel-up}{819.847}

% --- The split (rebuilt as hgb/link_data.py builds it; checked against the fingerprint) ----------
\setval{split-train}{655.879}
\setval{split-val}{81.984}
\setval{split-test}{81.984}
\setval{split-seed}{1}
\setval{split-fingerprint}{c5f7bccbf61c1b0f}
\setval{split-negatives}{99}
\setval{train-deg-med}{6}  % training edges per user
\setval{train-deg-mean}{12,2}
\setval{train-deg-max}{1.745}
\setval{test-users}{33.222}  % users with at least one test edge
\setval{group-book}{23.754}
\setval{group-book-pct}{48,6}
\setval{test-book-pct}{41,0}  % share of test edges
\setval{group-dvd}{4.359}
\setval{group-dvd-pct}{8,9}
\setval{test-dvd-pct}{13,9}  % share of test edges
\setval{group-music}{16.436}
\setval{group-music-pct}{33,6}
\setval{test-music-pct}{33,8}  % share of test edges
\setval{group-video}{4.334}
\setval{group-video-pct}{8,9}
\setval{test-video-pct}{11,3}  % share of test edges
\setval{group-bookmusic-pct}{82,2}
\setval{test-bookmusic-pct}{74,8}  % books and music, share of test edges

% --- The graph: category links, near-empty paths and edition twins (audits, link.dat, HAN logs) ---
\setval{before-topuser-pct}{61}  % share of products the largest user had reviewed before the fix
\setval{raw-shared-pct}{68,7}  % raw review lines that also appear under another product
\setval{cat-all-median}{13}  % categories per product if every ancestor were linked
\setval{cat-all-largest-pct}{48,7}  % products in one category, every ancestor linked
\setval{cat-pcp-ratio}{20,7}  % P-C-P walks with every ancestor linked, over leaves only
\setval{pcp-isolated}{9}  % products whose leaf categories hold no other product
\setval{pccp-reach-train-pct}{27,7}  % training edges whose product PCCP links to another product
\setval{pccp-isolated-pct}{81,1}  % products PCCP links to no other product (full graph)
\setval{pccp-reach-book-pct}{2,2}  % share of Book products that PCCP links to another product
\setval{pccp-reach-dvd-pct}{90,1}  % share of DVD products that PCCP links to another product
\setval{pccp-reach-music-pct}{5,3}  % share of Music products that PCCP links to another product
\setval{pccp-reach-video-pct}{90,2}  % share of Video products that PCCP links to another product
\setval{empty-isolated-min-pct}{76,2}  % near-empty paths: products without a neighbour, HAN graphs
\setval{empty-isolated-max-pct}{81,1}
\setval{pccp-graph-isolated-pct}{81,1}  % PCCP in the HAN graph (pruning adds a few)
\setval{pcccp-graph-isolated}{0}
\setval{pcp-graph-isolated}{9}
\setval{pp-isolated-pct}{11,3}  % products without an also-bought partner
\setval{twin-products}{51}  % products with an identical reviewer set elsewhere
\setval{twin-groups}{23}
\setval{twin-edges}{497}
\setval{twin-val}{47}
\setval{twin-test}{40}
\setval{example-shared}{659}  % review lines a DVD and a VHS of one film share (review part 1)
\setval{example-lines}{660}

% --- The search: 30 runs per set size (LMSPS notes, Parts 1 and 5-9) -----------------------------
\setval{search-cands}{36}
\setval{search-pool}{10}
\setval{search-sets}{210}  % sets scored per run (10 choose 4 = 10 choose 6)
\setval{search-epochs}{150}
\setval{search-warmup}{15}
\setval{search-hops}{5}
\setval{search-hidden}{512}
\setval{search-batch}{10.000}
\setval{search-dropout}{0,5}
\setval{search-lr}{0,001}
\setval{search-tau-max}{8}
\setval{search-tau-min}{4}
\setval{search-delta}{0,010}
\setval{search-folds}{4}  % folds of training edges for the weight step
\setval{search-alr}{0,001}  % learning rate of the alpha step
\setval{search-weight-decay}{0}
\setval{search-input-drop}{0}
\setval{env-gpu}{RTX 4070 SUPER}  % the run records
\setval{env-gpu-mem}{12}  % GB, the card's specification
\setval{env-python}{3.9.25}
\setval{env-torch}{2.1.2}
\setval{env-cuda}{12.1}
\setval{env-dgl}{2.2.1}
\setval{env-numpy}{1.26.4}
\setval{search-runs}{30}
\setval{search-runs-total}{60}
\setval{search-cands-raw}{55}  % P...P paths of at most 5 hops before merging reversed pairs
\setval{search-epoch-s-4}{6,9}  % median seconds per training epoch
\setval{search-prop-s-4}{37}  % median seconds of feature propagation
\setval{search-run-min-4}{18}  % median minutes per run (epochs + propagation)
\setval{search-fleet-h-4}{9,0}  % GPU hours for the 30 runs
\setval{search-winners-4}{26}  % distinct winners among the 30 runs
\setval{search-delta-alt}{0,005}  % the second tolerance of the vote (sensitivity check)
\setval{search-tie-med-4}{60}  % median tie-group size (delta 0.010)
\setval{search-tie-min-4}{13}
\setval{search-tie-max-4}{87}
\setval{search-tie-seed1-4}{83}  % tie group of seed 1
\setval{search-noise-4}{0,0055}  % score noise of one set
\setval{search-agree-4}{\ensuremath{+0,68}}  % Spearman of two runs on shared sets
\setval{search-xnoise-4}{0,0050}  % gap of a winner re-scored in another run
\setval{search-vote-4}{16}  % distinct vote picks
\setval{search-null-med-4}{21}
\setval{search-null-lo-4}{17}  % 90% range, as reselect_fleet.py
\setval{search-null-hi-4}{26}
\setval{search-null-share-4}{3,0}  % null fleets with as few distinct picks
\setval{search-vote-alt-4}{20}  % distinct vote picks
\setval{search-null-med-alt-4}{26}
\setval{search-null-lo-alt-4}{23}  % 90% range, as reselect_fleet.py
\setval{search-null-hi-alt-4}{29}
\setval{search-null-share-alt-4}{0,0}  % null fleets with as few distinct picks
\setval{search-cons-count-4}{5}  % runs whose vote pick is the consensus
\setval{search-null-fleets}{200}  % shuffled-score fleets in the null
\setval{search-avg-sets-4}{333}  % sets scored in at least 5 runs
\setval{search-se-4}{0,0025}  % standard error of the best averaged set
\setval{search-window-4}{2}  % sets within one SE of the best averaged score
\setval{search-epoch-s-6}{8,1}  % median seconds per training epoch
\setval{search-prop-s-6}{36}  % median seconds of feature propagation
\setval{search-run-min-6}{21}  % median minutes per run (epochs + propagation)
\setval{search-fleet-h-6}{10,7}  % GPU hours for the 30 runs
\setval{search-winners-6}{30}  % distinct winners among the 30 runs
\setval{search-tie-med-6}{78}  % median tie-group size (delta 0.010)
\setval{search-tie-min-6}{50}
\setval{search-tie-max-6}{146}
\setval{search-tie-seed1-6}{72}  % tie group of seed 1
\setval{search-noise-6}{0,0068}  % score noise of one set
\setval{search-agree-6}{\ensuremath{+0,59}}  % Spearman of two runs on shared sets
\setval{search-xnoise-6}{0,0039}  % gap of a winner re-scored in another run
\setval{search-vote-6}{19}  % distinct vote picks
\setval{search-null-med-6}{22}
\setval{search-null-lo-6}{19}  % 90% range, as reselect_fleet.py
\setval{search-null-hi-6}{25}
\setval{search-null-share-6}{11,0}  % null fleets with as few distinct picks
\setval{search-vote-alt-6}{20}  % distinct vote picks
\setval{search-null-med-alt-6}{25}
\setval{search-null-lo-alt-6}{22}  % 90% range, as reselect_fleet.py
\setval{search-null-hi-alt-6}{28}
\setval{search-null-share-alt-6}{1,0}  % null fleets with as few distinct picks
\setval{search-cons-count-6}{5}  % runs whose vote pick is the consensus
\setval{search-avg-sets-6}{169}  % sets scored in at least 5 runs
\setval{search-se-6}{0,0030}  % standard error of the best averaged set
\setval{search-window-6}{32}  % sets within one SE of the best averaged score
\setval{pilot-sets}{210}  % sets scored in each same-seed pilot run, all identical
\setval{search-total-h}{20}  % GPU hours for all 60 runs (epochs + propagation)
\setval{rank-pupup-4}{3,6}  % mean alpha rank of PUPUP (1 = most preferred of 36)
\setval{rank-pupup-6}{3,6}  % mean alpha rank of PUPUP (1 = most preferred of 36)
\setval{rank-puppup-4}{2,8}  % mean alpha rank of PUPPUP (1 = most preferred of 36)
\setval{rank-puppup-6}{4,1}  % mean alpha rank of PUPPUP (1 = most preferred of 36)
\setval{rank-pup-4}{9,1}  % mean alpha rank of PUP (1 = most preferred of 36)
\setval{rank-pup-6}{10,2}  % mean alpha rank of PUP (1 = most preferred of 36)
\setval{rank-pcp-4}{26,6}  % mean alpha rank of PCP (1 = most preferred of 36)
\setval{rank-pcp-6}{28,7}  % mean alpha rank of PCP (1 = most preferred of 36)
\setval{rank-pcccp-4}{22,3}  % mean alpha rank of PCCCP (1 = most preferred of 36)
\setval{rank-pcccp-6}{21,8}  % mean alpha rank of PCCCP (1 = most preferred of 36)
\setval{rank-pcpcp-4}{24,5}  % mean alpha rank of PCPCP (1 = most preferred of 36)
\setval{rank-pcpcp-6}{27,5}  % mean alpha rank of PCPCP (1 = most preferred of 36)
\setval{rank-ppccpp-4}{14,9}  % mean alpha rank of PPCCPP (1 = most preferred of 36)
\setval{rank-ppccpp-6}{6,4}  % mean alpha rank of PPCCPP (1 = most preferred of 36)

% --- Settings: the commands in the seed-1 logs of M5 (HAN, metapath2vec) and of GAT --------------
\setval{han-set-topk}{20}  % --topk
\setval{han-set-hidden}{512}  % --hidden
\setval{han-set-num-heads}{8}  % --num-heads
\setval{han-set-feat-drop}{0,2}  % --feat-drop
\setval{han-set-attn-drop}{0,2}  % --attn-drop
\setval{han-set-lr}{0,001}  % --lr
\setval{han-set-weight-decay}{0}  % --weight-decay
\setval{han-set-batch-size}{100.000}  % --batch-size
\setval{han-set-stage}{150}  % --stage
\setval{han-set-patience}{10}  % --patience
\setval{han-set-min-delta}{0,0005}  % --min-delta
\setval{han-set-beta-warmup}{15}  % --beta-warmup
\setval{mtv-set-dim}{512}  % --dim
\setval{mtv-set-walks-per-node}{5}  % --walks-per-node
\setval{mtv-set-walk-length}{20}  % --walk-length
\setval{mtv-set-window}{3}  % --window
\setval{mtv-set-negatives}{5}  % --negatives
\setval{mtv-set-lr}{0,01}  % --lr
\setval{mtv-set-epochs}{10}  % --epochs
\setval{mtv-set-evals-per-epoch}{4}  % --evals-per-epoch
\setval{mtv-set-patience}{6}  % --patience
\setval{mtv-set-min-delta}{0,0005}  % --min-delta
\setval{mtv-set-batch-walks}{1.024}  % --batch-walks
\setval{gat-set-in-dim}{256}  % --in-dim
\setval{gat-set-hidden}{512}  % --hidden
\setval{gat-set-num-heads}{8}  % --num-heads
\setval{gat-set-layers}{2}  % --layers
\setval{gat-set-feat-drop}{0,2}  % --feat-drop
\setval{gat-set-attn-drop}{0,2}  % --attn-drop
\setval{gat-set-lr}{0,001}  % --lr
\setval{gat-set-weight-decay}{0}  % --weight-decay
\setval{gat-set-batch-size}{100.000}  % --batch-size
\setval{gat-set-stage}{150}  % --stage
\setval{gat-set-patience}{10}  % --patience
\setval{gat-set-min-delta}{0,0005}  % --min-delta
\setval{han-set-head-dim}{64}  % dimensions per attention head
\setval{han-set-sem-dim}{128}  % semantic attention vector (model print in the log)
\setval{han-set-in-dim}{260}  % product input features: LSA title vector and product group
\setval{han-set-semantic-tau}{1}  % softmax temperature

% --- Path sets (HAN repository, arch.py) ---------------------------------------------------------
\setval{runs-total}{68}  % result rows: every model, set and seed
\setval{sets-distinct-paths}{32}  % distinct paths over the 24 sets, all among the candidates
\setval{sets-size-min}{2}  % paths per set, over the 24 sets
\setval{sets-size-max}{6}
\setval{sets-hops-max}{5}  % longest path in the 24 sets, hops
\setval{perm-shuffles}{20.000}  % permutation tests of the Spearman correlations
\setval{set-m1}{PUP PCP}  % amazon_lp_manual_1
\setval{set-m1-n}{2}
\setval{set-m1-hops}{4}
\setval{set-m2}{PP PUP PCP}  % amazon_lp_manual_2
\setval{set-m2-n}{3}
\setval{set-m2-hops}{5}
\setval{set-m3}{PP PUP PPUP}  % amazon_lp_manual_3
\setval{set-m3-n}{3}
\setval{set-m3-hops}{6}
\setval{set-m4}{PCP PCCCP PCPCP}  % amazon_lp_manual_4
\setval{set-m4-n}{3}
\setval{set-m4-hops}{10}
\setval{set-m5}{PP PUP PCP PPUP}  % amazon_lp_manual_5
\setval{set-m5-n}{4}
\setval{set-m5-hops}{8}
\setval{set-m6}{PP PUP PCP PPUP PPP PCCCP}  % amazon_lp_manual_6
\setval{set-m6-n}{6}
\setval{set-m6-hops}{14}
\setval{set-m7}{PP PPP PUP PCP}  % amazon_lp_manual_7
\setval{set-m7-n}{4}
\setval{set-m7-hops}{7}
\setval{set-m8}{PUP PCP PUPUP PCPUP}  % amazon_lp_manual_8
\setval{set-m8-n}{4}
\setval{set-m8-hops}{12}
\setval{set-r4d1}{PCCPPP PCPPUP PPCCPP PPPPP}  % amazon_lp_rand4_d1
\setval{set-r4d1-n}{4}
\setval{set-r4d1-hops}{19}
\setval{set-r4d2}{PCCPCP PCCPP PCPCP PPUP}  % amazon_lp_rand4_d2
\setval{set-r4d2-n}{4}
\setval{set-r4d2-hops}{16}
\setval{set-r4d3}{PCCCP PP PPPPP PUPUP}  % amazon_lp_rand4_d3
\setval{set-r4d3-n}{4}
\setval{set-r4d3-hops}{13}
\setval{set-r4d4}{PCCP PCPPPP PUP PUPPUP}  % amazon_lp_rand4_d4
\setval{set-r4d4-n}{4}
\setval{set-r4d4-hops}{15}
\setval{set-r4d5}{PCCPPP PCPP PCPUPP PPP}  % amazon_lp_rand4_d5
\setval{set-r4d5-n}{4}
\setval{set-r4d5-hops}{15}
\setval{set-r6d1}{PCCCCP PCCCPP PCCPCP PCCPP PCPUP PPUPP}  % amazon_lp_rand6_d1
\setval{set-r6d1-n}{6}
\setval{set-r6d1-hops}{27}
\setval{set-r6d2}{PCCCP PCPPUP PCPUPP PPCPP PPP PPUP}  % amazon_lp_rand6_d2
\setval{set-r6d2-n}{6}
\setval{set-r6d2-hops}{23}
\setval{set-r6d3}{PCCPPP PCPPCP PCPUP PPCPPP PPPPPP PUPUP}  % amazon_lp_rand6_d3
\setval{set-r6d3-n}{6}
\setval{set-r6d3-hops}{28}
\setval{set-r6d4}{PCCPP PPCPPP PPP PPPPP PPPPUP PUPUP}  % amazon_lp_rand6_d4
\setval{set-r6d4-n}{6}
\setval{set-r6d4-hops}{24}
\setval{set-r6d5}{PCCPPP PCPPCP PCPUP PPCPUP PPPPPP PPPUP}  % amazon_lp_rand6_d5
\setval{set-r6d5-n}{6}
\setval{set-r6d5-hops}{28}
\setval{set-f4a1}{PCCPCP PPCCPP PPUPP PUPUP}  % amazon_lp_v3_f4_a1
\setval{set-f4a1-n}{4}
\setval{set-f4a1-hops}{18}
\setval{set-f4a2}{PCCPPP PPCCPP PPPUPP PUPUP}  % amazon_lp_v3_f4_a2
\setval{set-f4a2-n}{4}
\setval{set-f4a2-hops}{19}
\setval{set-f4c}{PCCPP PPUPP PUPPUP PUPUP}  % amazon_lp_v3_f4_c
\setval{set-f4c-n}{4}
\setval{set-f4c-hops}{17}
\setval{set-f6a1}{PCCP PCCPP PPCCPP PPUP PUPPUP PUPUP}  % amazon_lp_v3_f6_a1
\setval{set-f6a1-n}{6}
\setval{set-f6a1-hops}{24}
\setval{set-f6a2}{PCCP PCCPCP PPCCPP PPUP PUPPUP PUPUP}  % amazon_lp_v3_f6_a2
\setval{set-f6a2-n}{6}
\setval{set-f6a2-hops}{25}
\setval{set-f6c}{PCCP PCCPUP PPCCPP PPUPP PUPPUP PUPUP}  % amazon_lp_v3_f6_c
\setval{set-f6c-n}{6}
\setval{set-f6c-hops}{26}

% --- Every set at seed 1: HAN and metapath2vec (results.csv) -------------------------------------
\setval{han-m1-hit-s1}{0,635}
\setval{han-m1-ndcg-s1}{0,398}
\setval{han-m1-fhit-s1}{0,0169}
\setval{han-m1-fndcg-s1}{0,0082}
\setval{han-m1-frec-s1}{0,0305}
\setval{han-m1-val-s1}{0,3072}
\setval{han-m2-hit-s1}{0,650}
\setval{han-m2-ndcg-s1}{0,413}
\setval{han-m2-fhit-s1}{0,0192}
\setval{han-m2-fndcg-s1}{0,0093}
\setval{han-m2-frec-s1}{0,0337}
\setval{han-m2-val-s1}{0,2957}
\setval{han-m3-hit-s1}{0,648}
\setval{han-m3-ndcg-s1}{0,416}
\setval{han-m3-fhit-s1}{0,0203}
\setval{han-m3-fndcg-s1}{0,0102}
\setval{han-m3-frec-s1}{0,0382}
\setval{han-m3-val-s1}{0,2935}
\setval{han-m4-hit-s1}{0,619}
\setval{han-m4-ndcg-s1}{0,375}
\setval{han-m4-fhit-s1}{0,0121}
\setval{han-m4-fndcg-s1}{0,0058}
\setval{han-m4-frec-s1}{0,0244}
\setval{han-m4-val-s1}{0,3194}
\setval{han-m5-hit-s1}{0,665}
\setval{han-m5-ndcg-s1}{0,427}
\setval{han-m5-fhit-s1}{0,0216}
\setval{han-m5-fndcg-s1}{0,0106}
\setval{han-m5-frec-s1}{0,0388}
\setval{han-m5-val-s1}{0,2814}
\setval{han-m6-hit-s1}{0,692}
\setval{han-m6-ndcg-s1}{0,453}
\setval{han-m6-fhit-s1}{0,0257}
\setval{han-m6-fndcg-s1}{0,0128}
\setval{han-m6-frec-s1}{0,0462}
\setval{han-m6-val-s1}{0,2651}
\setval{han-m7-hit-s1}{0,670}
\setval{han-m7-ndcg-s1}{0,433}
\setval{han-m7-fhit-s1}{0,0228}
\setval{han-m7-fndcg-s1}{0,0114}
\setval{han-m7-frec-s1}{0,0423}
\setval{han-m7-val-s1}{0,2834}
\setval{han-m8-hit-s1}{0,663}
\setval{han-m8-ndcg-s1}{0,427}
\setval{han-m8-fhit-s1}{0,0212}
\setval{han-m8-fndcg-s1}{0,0102}
\setval{han-m8-frec-s1}{0,0385}
\setval{han-m8-val-s1}{0,2875}
\setval{han-r4d1-hit-s1}{0,671}
\setval{han-r4d1-ndcg-s1}{0,441}
\setval{han-r4d1-fhit-s1}{0,0267}
\setval{han-r4d1-fndcg-s1}{0,0133}
\setval{han-r4d1-frec-s1}{0,0501}
\setval{han-r4d1-val-s1}{0,2816}
\setval{han-r4d2-hit-s1}{0,662}
\setval{han-r4d2-ndcg-s1}{0,421}
\setval{han-r4d2-fhit-s1}{0,0203}
\setval{han-r4d2-fndcg-s1}{0,0102}
\setval{han-r4d2-frec-s1}{0,0381}
\setval{han-r4d2-val-s1}{0,2865}
\setval{han-r4d3-hit-s1}{0,677}
\setval{han-r4d3-ndcg-s1}{0,443}
\setval{han-r4d3-fhit-s1}{0,0243}
\setval{han-r4d3-fndcg-s1}{0,0124}
\setval{han-r4d3-frec-s1}{0,0459}
\setval{han-r4d3-val-s1}{0,2785}
\setval{han-r4d4-hit-s1}{0,671}
\setval{han-r4d4-ndcg-s1}{0,436}
\setval{han-r4d4-fhit-s1}{0,0227}
\setval{han-r4d4-fndcg-s1}{0,0113}
\setval{han-r4d4-frec-s1}{0,0413}
\setval{han-r4d4-val-s1}{0,2827}
\setval{han-r4d5-hit-s1}{0,660}
\setval{han-r4d5-ndcg-s1}{0,422}
\setval{han-r4d5-fhit-s1}{0,0229}
\setval{han-r4d5-fndcg-s1}{0,0113}
\setval{han-r4d5-frec-s1}{0,0435}
\setval{han-r4d5-val-s1}{0,2922}
\setval{han-r6d1-hit-s1}{0,673}
\setval{han-r6d1-ndcg-s1}{0,436}
\setval{han-r6d1-fhit-s1}{0,0228}
\setval{han-r6d1-fndcg-s1}{0,0113}
\setval{han-r6d1-frec-s1}{0,0403}
\setval{han-r6d1-val-s1}{0,2821}
\setval{han-r6d2-hit-s1}{0,706}
\setval{han-r6d2-ndcg-s1}{0,468}
\setval{han-r6d2-fhit-s1}{0,0280}
\setval{han-r6d2-fndcg-s1}{0,0142}
\setval{han-r6d2-frec-s1}{0,0501}
\setval{han-r6d2-val-s1}{0,2580}
\setval{han-r6d3-hit-s1}{0,694}
\setval{han-r6d3-ndcg-s1}{0,460}
\setval{han-r6d3-fhit-s1}{0,0271}
\setval{han-r6d3-fndcg-s1}{0,0135}
\setval{han-r6d3-frec-s1}{0,0503}
\setval{han-r6d3-val-s1}{0,2708}
\setval{han-r6d4-hit-s1}{0,675}
\setval{han-r6d4-ndcg-s1}{0,443}
\setval{han-r6d4-fhit-s1}{0,0262}
\setval{han-r6d4-fndcg-s1}{0,0133}
\setval{han-r6d4-frec-s1}{0,0485}
\setval{han-r6d4-val-s1}{0,2782}
\setval{han-r6d5-hit-s1}{0,706}
\setval{han-r6d5-ndcg-s1}{0,469}
\setval{han-r6d5-fhit-s1}{0,0282}
\setval{han-r6d5-fndcg-s1}{0,0142}
\setval{han-r6d5-frec-s1}{0,0522}
\setval{han-r6d5-val-s1}{0,2610}
\setval{han-f4a1-hit-s1}{0,635}
\setval{han-f4a1-ndcg-s1}{0,405}
\setval{han-f4a1-fhit-s1}{0,0205}
\setval{han-f4a1-fndcg-s1}{0,0104}
\setval{han-f4a1-frec-s1}{0,0368}
\setval{han-f4a1-val-s1}{0,3056}
\setval{han-f4a2-hit-s1}{0,650}
\setval{han-f4a2-ndcg-s1}{0,421}
\setval{han-f4a2-fhit-s1}{0,0232}
\setval{han-f4a2-fndcg-s1}{0,0116}
\setval{han-f4a2-frec-s1}{0,0407}
\setval{han-f4a2-val-s1}{0,3019}
\setval{han-f4c-hit-s1}{0,658}
\setval{han-f4c-ndcg-s1}{0,430}
\setval{han-f4c-fhit-s1}{0,0233}
\setval{han-f4c-fndcg-s1}{0,0115}
\setval{han-f4c-frec-s1}{0,0441}
\setval{han-f4c-val-s1}{0,2984}
\setval{han-f6a1-hit-s1}{0,666}
\setval{han-f6a1-ndcg-s1}{0,435}
\setval{han-f6a1-fhit-s1}{0,0246}
\setval{han-f6a1-fndcg-s1}{0,0124}
\setval{han-f6a1-frec-s1}{0,0449}
\setval{han-f6a1-val-s1}{0,2885}
\setval{han-f6a2-hit-s1}{0,666}
\setval{han-f6a2-ndcg-s1}{0,435}
\setval{han-f6a2-fhit-s1}{0,0238}
\setval{han-f6a2-fndcg-s1}{0,0121}
\setval{han-f6a2-frec-s1}{0,0448}
\setval{han-f6a2-val-s1}{0,2887}
\setval{han-f6c-hit-s1}{0,666}
\setval{han-f6c-ndcg-s1}{0,435}
\setval{han-f6c-fhit-s1}{0,0249}
\setval{han-f6c-fndcg-s1}{0,0127}
\setval{han-f6c-frec-s1}{0,0451}
\setval{han-f6c-val-s1}{0,2885}
\setval{mtv-m1-hit-s1}{0,600}
\setval{mtv-m1-ndcg-s1}{0,374}
\setval{mtv-m1-fhit-s1}{0,0122}
\setval{mtv-m1-fndcg-s1}{0,0061}
\setval{mtv-m1-frec-s1}{0,0188}
\setval{mtv-m1-val-s1}{0,3401}
\setval{mtv-m2-hit-s1}{0,638}
\setval{mtv-m2-ndcg-s1}{0,424}
\setval{mtv-m2-fhit-s1}{0,0263}
\setval{mtv-m2-fndcg-s1}{0,0129}
\setval{mtv-m2-frec-s1}{0,0459}
\setval{mtv-m2-val-s1}{0,3174}
\setval{mtv-m3-hit-s1}{0,624}
\setval{mtv-m3-ndcg-s1}{0,431}
\setval{mtv-m3-fhit-s1}{0,0446}
\setval{mtv-m3-fndcg-s1}{0,0230}
\setval{mtv-m3-frec-s1}{0,0803}
\setval{mtv-m3-val-s1}{0,3260}
\setval{mtv-m4-hit-s1}{0,095}
\setval{mtv-m4-ndcg-s1}{0,042}
\setval{mtv-m4-fhit-s1}{0,0002}
\setval{mtv-m4-fndcg-s1}{0,0001}
\setval{mtv-m4-frec-s1}{0,0003}
\setval{mtv-m4-val-s1}{0,6931}
\setval{mtv-m5-hit-s1}{0,639}
\setval{mtv-m5-ndcg-s1}{0,427}
\setval{mtv-m5-fhit-s1}{0,0273}
\setval{mtv-m5-fndcg-s1}{0,0136}
\setval{mtv-m5-frec-s1}{0,0452}
\setval{mtv-m5-val-s1}{0,3119}
\setval{mtv-m6-hit-s1}{0,639}
\setval{mtv-m6-ndcg-s1}{0,432}
\setval{mtv-m6-fhit-s1}{0,0315}
\setval{mtv-m6-fndcg-s1}{0,0157}
\setval{mtv-m6-frec-s1}{0,0568}
\setval{mtv-m6-val-s1}{0,3166}
\setval{mtv-m7-hit-s1}{0,633}
\setval{mtv-m7-ndcg-s1}{0,421}
\setval{mtv-m7-fhit-s1}{0,0272}
\setval{mtv-m7-fndcg-s1}{0,0135}
\setval{mtv-m7-frec-s1}{0,0498}
\setval{mtv-m7-val-s1}{0,3188}
\setval{mtv-m8-hit-s1}{0,619}
\setval{mtv-m8-ndcg-s1}{0,392}
\setval{mtv-m8-fhit-s1}{0,0153}
\setval{mtv-m8-fndcg-s1}{0,0075}
\setval{mtv-m8-frec-s1}{0,0248}
\setval{mtv-m8-val-s1}{0,3266}
\setval{mtv-r4d1-hit-s1}{0,553}
\setval{mtv-r4d1-ndcg-s1}{0,360}
\setval{mtv-r4d1-fhit-s1}{0,0179}
\setval{mtv-r4d1-fndcg-s1}{0,0090}
\setval{mtv-r4d1-frec-s1}{0,0353}
\setval{mtv-r4d1-val-s1}{0,4152}
\setval{mtv-r4d2-hit-s1}{0,583}
\setval{mtv-r4d2-ndcg-s1}{0,368}
\setval{mtv-r4d2-fhit-s1}{0,0107}
\setval{mtv-r4d2-fndcg-s1}{0,0050}
\setval{mtv-r4d2-frec-s1}{0,0168}
\setval{mtv-r4d2-val-s1}{0,3733}
\setval{mtv-r4d3-hit-s1}{0,607}
\setval{mtv-r4d3-ndcg-s1}{0,406}
\setval{mtv-r4d3-fhit-s1}{0,0316}
\setval{mtv-r4d3-fndcg-s1}{0,0158}
\setval{mtv-r4d3-frec-s1}{0,0589}
\setval{mtv-r4d3-val-s1}{0,3378}
\setval{mtv-r4d4-hit-s1}{0,631}
\setval{mtv-r4d4-ndcg-s1}{0,419}
\setval{mtv-r4d4-fhit-s1}{0,0263}
\setval{mtv-r4d4-fndcg-s1}{0,0136}
\setval{mtv-r4d4-frec-s1}{0,0436}
\setval{mtv-r4d4-val-s1}{0,3151}
\setval{mtv-r4d5-hit-s1}{0,548}
\setval{mtv-r4d5-ndcg-s1}{0,348}
\setval{mtv-r4d5-fhit-s1}{0,0149}
\setval{mtv-r4d5-fndcg-s1}{0,0074}
\setval{mtv-r4d5-frec-s1}{0,0280}
\setval{mtv-r4d5-val-s1}{0,4075}
\setval{mtv-r6d1-hit-s1}{0,583}
\setval{mtv-r6d1-ndcg-s1}{0,368}
\setval{mtv-r6d1-fhit-s1}{0,0128}
\setval{mtv-r6d1-fndcg-s1}{0,0061}
\setval{mtv-r6d1-frec-s1}{0,0230}
\setval{mtv-r6d1-val-s1}{0,3762}
\setval{mtv-r6d2-hit-s1}{0,622}
\setval{mtv-r6d2-ndcg-s1}{0,409}
\setval{mtv-r6d2-fhit-s1}{0,0226}
\setval{mtv-r6d2-fndcg-s1}{0,0110}
\setval{mtv-r6d2-frec-s1}{0,0393}
\setval{mtv-r6d2-val-s1}{0,3374}
\setval{mtv-r6d3-hit-s1}{0,623}
\setval{mtv-r6d3-ndcg-s1}{0,403}
\setval{mtv-r6d3-fhit-s1}{0,0205}
\setval{mtv-r6d3-fndcg-s1}{0,0100}
\setval{mtv-r6d3-frec-s1}{0,0373}
\setval{mtv-r6d3-val-s1}{0,3237}
\setval{mtv-r6d4-hit-s1}{0,632}
\setval{mtv-r6d4-ndcg-s1}{0,436}
\setval{mtv-r6d4-fhit-s1}{0,0391}
\setval{mtv-r6d4-fndcg-s1}{0,0198}
\setval{mtv-r6d4-frec-s1}{0,0761}
\setval{mtv-r6d4-val-s1}{0,3190}
\setval{mtv-r6d5-hit-s1}{0,607}
\setval{mtv-r6d5-ndcg-s1}{0,394}
\setval{mtv-r6d5-fhit-s1}{0,0187}
\setval{mtv-r6d5-fndcg-s1}{0,0091}
\setval{mtv-r6d5-frec-s1}{0,0340}
\setval{mtv-r6d5-val-s1}{0,3509}
\setval{mtv-f4a1-hit-s1}{0,605}
\setval{mtv-f4a1-ndcg-s1}{0,400}
\setval{mtv-f4a1-fhit-s1}{0,0290}
\setval{mtv-f4a1-fndcg-s1}{0,0150}
\setval{mtv-f4a1-frec-s1}{0,0473}
\setval{mtv-f4a1-val-s1}{0,3442}
\setval{mtv-f4a2-hit-s1}{0,608}
\setval{mtv-f4a2-ndcg-s1}{0,405}
\setval{mtv-f4a2-fhit-s1}{0,0302}
\setval{mtv-f4a2-fndcg-s1}{0,0158}
\setval{mtv-f4a2-frec-s1}{0,0516}
\setval{mtv-f4a2-val-s1}{0,3409}
\setval{mtv-f4c-hit-s1}{0,626}
\setval{mtv-f4c-ndcg-s1}{0,421}
\setval{mtv-f4c-fhit-s1}{0,0364}
\setval{mtv-f4c-fndcg-s1}{0,0193}
\setval{mtv-f4c-frec-s1}{0,0582}
\setval{mtv-f4c-val-s1}{0,3222}
\setval{mtv-f6a1-hit-s1}{0,588}
\setval{mtv-f6a1-ndcg-s1}{0,385}
\setval{mtv-f6a1-fhit-s1}{0,0265}
\setval{mtv-f6a1-fndcg-s1}{0,0142}
\setval{mtv-f6a1-frec-s1}{0,0399}
\setval{mtv-f6a1-val-s1}{0,3528}
\setval{mtv-f6a2-hit-s1}{0,592}
\setval{mtv-f6a2-ndcg-s1}{0,384}
\setval{mtv-f6a2-fhit-s1}{0,0252}
\setval{mtv-f6a2-fndcg-s1}{0,0132}
\setval{mtv-f6a2-frec-s1}{0,0381}
\setval{mtv-f6a2-val-s1}{0,3503}
\setval{mtv-f6c-hit-s1}{0,598}
\setval{mtv-f6c-ndcg-s1}{0,392}
\setval{mtv-f6c-fhit-s1}{0,0266}
\setval{mtv-f6c-fndcg-s1}{0,0141}
\setval{mtv-f6c-frec-s1}{0,0403}
\setval{mtv-f6c-val-s1}{0,3435}

% --- The main comparison: 3 seeds, Welch tests (HAN notes, Part 10) ------------------------------
\setval{han-f4a1-hit-mean}{0,635}
\setval{han-f4a1-hit-sd}{0,003}
\setval{han-m5-hit-mean}{0,664}
\setval{han-m5-hit-sd}{0,002}
\setval{han-4-hit-diff}{\ensuremath{-0,029}}  % f4a1 minus m5
\setval{han-4-hit-rel}{\ensuremath{-4,3}}  % relative difference, %
\setval{han-4-hit-p}{\ensuremath{<}0,001}  % Welch, two-sided
\setval{han-4-hit-pt}{\ensuremath{p<0,001}}
\setval{han-4-hit-diffabs}{0,029}
\setval{han-4-hit-relabs}{4,3}
\setval{han-4-hit-sdx}{10}  % difference in units of the larger seed SD
\setval{han-f4a1-ndcg-mean}{0,405}
\setval{han-f4a1-ndcg-sd}{0,002}
\setval{han-m5-ndcg-mean}{0,426}
\setval{han-m5-ndcg-sd}{0,002}
\setval{han-4-ndcg-diff}{\ensuremath{-0,020}}  % f4a1 minus m5
\setval{han-4-ndcg-rel}{\ensuremath{-4,8}}  % relative difference, %
\setval{han-4-ndcg-p}{\ensuremath{<}0,001}  % Welch, two-sided
\setval{han-4-ndcg-pt}{\ensuremath{p<0,001}}
\setval{han-4-ndcg-diffabs}{0,020}
\setval{han-4-ndcg-relabs}{4,8}
\setval{han-4-ndcg-sdx}{8}  % difference in units of the larger seed SD
\setval{han-f4a1-fhit-mean}{0,0207}
\setval{han-f4a1-fhit-sd}{0,0005}
\setval{han-m5-fhit-mean}{0,0211}
\setval{han-m5-fhit-sd}{0,0004}
\setval{han-4-fhit-diff}{\ensuremath{-0,0004}}  % f4a1 minus m5
\setval{han-4-fhit-rel}{\ensuremath{-2,0}}  % relative difference, %
\setval{han-4-fhit-p}{0,32}  % Welch, two-sided
\setval{han-4-fhit-pt}{\ensuremath{p=0,32}}
\setval{han-4-fhit-diffabs}{0,0004}
\setval{han-4-fhit-relabs}{2,0}
\setval{han-4-fhit-sdx}{1}  % difference in units of the larger seed SD
\setval{han-f4a1-fndcg-mean}{0,0103}
\setval{han-f4a1-fndcg-sd}{0,0002}
\setval{han-m5-fndcg-mean}{0,0104}
\setval{han-m5-fndcg-sd}{0,0002}
\setval{han-4-fndcg-diff}{\ensuremath{-0,0002}}  % f4a1 minus m5
\setval{han-4-fndcg-rel}{\ensuremath{-1,7}}  % relative difference, %
\setval{han-4-fndcg-p}{0,37}  % Welch, two-sided
\setval{han-4-fndcg-pt}{\ensuremath{p=0,37}}
\setval{han-4-fndcg-diffabs}{0,0002}
\setval{han-4-fndcg-relabs}{1,7}
\setval{han-4-fndcg-sdx}{1}  % difference in units of the larger seed SD
\setval{han-f4a1-frec-mean}{0,0376}
\setval{han-f4a1-frec-sd}{0,0007}
\setval{han-m5-frec-mean}{0,0379}
\setval{han-m5-frec-sd}{0,0008}
\setval{han-4-frec-diff}{\ensuremath{-0,0004}}  % f4a1 minus m5
\setval{han-4-frec-rel}{\ensuremath{-1,0}}  % relative difference, %
\setval{han-4-frec-p}{0,57}  % Welch, two-sided
\setval{han-4-frec-pt}{\ensuremath{p=0,57}}
\setval{han-4-frec-diffabs}{0,0004}
\setval{han-4-frec-relabs}{1,0}
\setval{han-4-frec-sdx}{0}  % difference in units of the larger seed SD
\setval{han-f6a1-hit-mean}{0,666}
\setval{han-f6a1-hit-sd}{0,002}
\setval{han-m6-hit-mean}{0,687}
\setval{han-m6-hit-sd}{0,004}
\setval{han-6-hit-diff}{\ensuremath{-0,021}}  % f6a1 minus m6
\setval{han-6-hit-rel}{\ensuremath{-3,0}}  % relative difference, %
\setval{han-6-hit-p}{0,0076}  % Welch, two-sided
\setval{han-6-hit-pt}{\ensuremath{p=0,0076}}
\setval{han-6-hit-diffabs}{0,021}
\setval{han-6-hit-relabs}{3,0}
\setval{han-6-hit-sdx}{5}  % difference in units of the larger seed SD
\setval{han-f6a1-ndcg-mean}{0,435}
\setval{han-f6a1-ndcg-sd}{0,001}
\setval{han-m6-ndcg-mean}{0,449}
\setval{han-m6-ndcg-sd}{0,004}
\setval{han-6-ndcg-diff}{\ensuremath{-0,014}}  % f6a1 minus m6
\setval{han-6-ndcg-rel}{\ensuremath{-3,1}}  % relative difference, %
\setval{han-6-ndcg-p}{0,017}  % Welch, two-sided
\setval{han-6-ndcg-pt}{\ensuremath{p=0,017}}
\setval{han-6-ndcg-diffabs}{0,014}
\setval{han-6-ndcg-relabs}{3,1}
\setval{han-6-ndcg-sdx}{3}  % difference in units of the larger seed SD
\setval{han-f6a1-fhit-mean}{0,0245}
\setval{han-f6a1-fhit-sd}{0,0004}
\setval{han-m6-fhit-mean}{0,0246}
\setval{han-m6-fhit-sd}{0,0011}
\setval{han-6-fhit-diff}{\ensuremath{-0,0001}}  % f6a1 minus m6
\setval{han-6-fhit-rel}{\ensuremath{-0,5}}  % relative difference, %
\setval{han-6-fhit-p}{0,87}  % Welch, two-sided
\setval{han-6-fhit-pt}{\ensuremath{p=0,87}}
\setval{han-6-fhit-diffabs}{0,0001}
\setval{han-6-fhit-relabs}{0,5}
\setval{han-6-fhit-sdx}{0}  % difference in units of the larger seed SD
\setval{han-f6a1-fndcg-mean}{0,0123}
\setval{han-f6a1-fndcg-sd}{0,0001}
\setval{han-m6-fndcg-mean}{0,0122}
\setval{han-m6-fndcg-sd}{0,0007}
\setval{han-6-fndcg-diff}{\ensuremath{+0,0001}}  % f6a1 minus m6
\setval{han-6-fndcg-rel}{\ensuremath{+0,6}}  % relative difference, %
\setval{han-6-fndcg-p}{0,87}  % Welch, two-sided
\setval{han-6-fndcg-pt}{\ensuremath{p=0,87}}
\setval{han-6-fndcg-diffabs}{0,0001}
\setval{han-6-fndcg-relabs}{0,6}
\setval{han-6-fndcg-sdx}{0}  % difference in units of the larger seed SD
\setval{han-f6a1-frec-mean}{0,0444}
\setval{han-f6a1-frec-sd}{0,0014}
\setval{han-m6-frec-mean}{0,0442}
\setval{han-m6-frec-sd}{0,0018}
\setval{han-6-frec-diff}{\ensuremath{+0,0002}}  % f6a1 minus m6
\setval{han-6-frec-rel}{\ensuremath{+0,5}}  % relative difference, %
\setval{han-6-frec-p}{0,88}  % Welch, two-sided
\setval{han-6-frec-pt}{\ensuremath{p=0,88}}
\setval{han-6-frec-diffabs}{0,0002}
\setval{han-6-frec-relabs}{0,5}
\setval{han-6-frec-sdx}{0}  % difference in units of the larger seed SD
\setval{mtv-f4a1-hit-mean}{0,600}
\setval{mtv-f4a1-hit-sd}{0,004}
\setval{mtv-m5-hit-mean}{0,640}
\setval{mtv-m5-hit-sd}{0,002}
\setval{mtv-4-hit-diff}{\ensuremath{-0,040}}  % f4a1 minus m5
\setval{mtv-4-hit-rel}{\ensuremath{-6,3}}  % relative difference, %
\setval{mtv-4-hit-p}{0,0013}  % Welch, two-sided
\setval{mtv-4-hit-pt}{\ensuremath{p=0,0013}}
\setval{mtv-4-hit-diffabs}{0,040}
\setval{mtv-4-hit-relabs}{6,3}
\setval{mtv-4-hit-sdx}{9}  % difference in units of the larger seed SD
\setval{mtv-f4a1-ndcg-mean}{0,395}
\setval{mtv-f4a1-ndcg-sd}{0,005}
\setval{mtv-m5-ndcg-mean}{0,426}
\setval{mtv-m5-ndcg-sd}{0,001}
\setval{mtv-4-ndcg-diff}{\ensuremath{-0,031}}  % f4a1 minus m5
\setval{mtv-4-ndcg-rel}{\ensuremath{-7,4}}  % relative difference, %
\setval{mtv-4-ndcg-p}{0,0071}  % Welch, two-sided
\setval{mtv-4-ndcg-pt}{\ensuremath{p=0,0071}}
\setval{mtv-4-ndcg-diffabs}{0,031}
\setval{mtv-4-ndcg-relabs}{7,4}
\setval{mtv-4-ndcg-sdx}{6}  % difference in units of the larger seed SD
\setval{mtv-f4a1-fhit-mean}{0,0279}
\setval{mtv-f4a1-fhit-sd}{0,0010}
\setval{mtv-m5-fhit-mean}{0,0272}
\setval{mtv-m5-fhit-sd}{0,0002}
\setval{mtv-4-fhit-diff}{\ensuremath{+0,0007}}  % f4a1 minus m5
\setval{mtv-4-fhit-rel}{\ensuremath{+2,5}}  % relative difference, %
\setval{mtv-4-fhit-p}{0,38}  % Welch, two-sided
\setval{mtv-4-fhit-pt}{\ensuremath{p=0,38}}
\setval{mtv-4-fhit-diffabs}{0,0007}
\setval{mtv-4-fhit-relabs}{2,5}
\setval{mtv-4-fhit-sdx}{1}  % difference in units of the larger seed SD
\setval{mtv-f4a1-fndcg-mean}{0,0147}
\setval{mtv-f4a1-fndcg-sd}{0,0004}
\setval{mtv-m5-fndcg-mean}{0,0136}
\setval{mtv-m5-fndcg-sd}{0,0002}
\setval{mtv-4-fndcg-diff}{\ensuremath{+0,0010}}  % f4a1 minus m5
\setval{mtv-4-fndcg-rel}{\ensuremath{+7,6}}  % relative difference, %
\setval{mtv-4-fndcg-p}{0,039}  % Welch, two-sided
\setval{mtv-4-fndcg-pt}{\ensuremath{p=0,039}}
\setval{mtv-4-fndcg-diffabs}{0,0010}
\setval{mtv-4-fndcg-relabs}{7,6}
\setval{mtv-4-fndcg-sdx}{2}  % difference in units of the larger seed SD
\setval{mtv-f4a1-frec-mean}{0,0434}
\setval{mtv-f4a1-frec-sd}{0,0034}
\setval{mtv-m5-frec-mean}{0,0452}
\setval{mtv-m5-frec-sd}{0,0005}
\setval{mtv-4-frec-diff}{\ensuremath{-0,0018}}  % f4a1 minus m5
\setval{mtv-4-frec-rel}{\ensuremath{-3,9}}  % relative difference, %
\setval{mtv-4-frec-p}{0,46}  % Welch, two-sided
\setval{mtv-4-frec-pt}{\ensuremath{p=0,46}}
\setval{mtv-4-frec-diffabs}{0,0018}
\setval{mtv-4-frec-relabs}{3,9}
\setval{mtv-4-frec-sdx}{1}  % difference in units of the larger seed SD
\setval{mtv-f6a1-hit-mean}{0,589}
\setval{mtv-f6a1-hit-sd}{0,002}
\setval{mtv-m6-hit-mean}{0,637}
\setval{mtv-m6-hit-sd}{0,002}
\setval{mtv-6-hit-diff}{\ensuremath{-0,049}}  % f6a1 minus m6
\setval{mtv-6-hit-rel}{\ensuremath{-7,6}}  % relative difference, %
\setval{mtv-6-hit-p}{\ensuremath{<}0,001}  % Welch, two-sided
\setval{mtv-6-hit-pt}{\ensuremath{p<0,001}}
\setval{mtv-6-hit-diffabs}{0,049}
\setval{mtv-6-hit-relabs}{7,6}
\setval{mtv-6-hit-sdx}{20}  % difference in units of the larger seed SD
\setval{mtv-f6a1-ndcg-mean}{0,384}
\setval{mtv-f6a1-ndcg-sd}{0,002}
\setval{mtv-m6-ndcg-mean}{0,431}
\setval{mtv-m6-ndcg-sd}{0,001}
\setval{mtv-6-ndcg-diff}{\ensuremath{-0,047}}  % f6a1 minus m6
\setval{mtv-6-ndcg-rel}{\ensuremath{-10,9}}  % relative difference, %
\setval{mtv-6-ndcg-p}{\ensuremath{<}0,001}  % Welch, two-sided
\setval{mtv-6-ndcg-pt}{\ensuremath{p<0,001}}
\setval{mtv-6-ndcg-diffabs}{0,047}
\setval{mtv-6-ndcg-relabs}{10,9}
\setval{mtv-6-ndcg-sdx}{21}  % difference in units of the larger seed SD
\setval{mtv-f6a1-fhit-mean}{0,0263}
\setval{mtv-f6a1-fhit-sd}{0,0005}
\setval{mtv-m6-fhit-mean}{0,0314}
\setval{mtv-m6-fhit-sd}{0,0006}
\setval{mtv-6-fhit-diff}{\ensuremath{-0,0051}}  % f6a1 minus m6
\setval{mtv-6-fhit-rel}{\ensuremath{-16,3}}  % relative difference, %
\setval{mtv-6-fhit-p}{\ensuremath{<}0,001}  % Welch, two-sided
\setval{mtv-6-fhit-pt}{\ensuremath{p<0,001}}
\setval{mtv-6-fhit-diffabs}{0,0051}
\setval{mtv-6-fhit-relabs}{16,3}
\setval{mtv-6-fhit-sdx}{9}  % difference in units of the larger seed SD
\setval{mtv-f6a1-fndcg-mean}{0,0140}
\setval{mtv-f6a1-fndcg-sd}{0,0004}
\setval{mtv-m6-fndcg-mean}{0,0157}
\setval{mtv-m6-fndcg-sd}{0,0001}
\setval{mtv-6-fndcg-diff}{\ensuremath{-0,0018}}  % f6a1 minus m6
\setval{mtv-6-fndcg-rel}{\ensuremath{-11,1}}  % relative difference, %
\setval{mtv-6-fndcg-p}{0,0084}  % Welch, two-sided
\setval{mtv-6-fndcg-pt}{\ensuremath{p=0,0084}}
\setval{mtv-6-fndcg-diffabs}{0,0018}
\setval{mtv-6-fndcg-relabs}{11,1}
\setval{mtv-6-fndcg-sdx}{5}  % difference in units of the larger seed SD
\setval{mtv-f6a1-frec-mean}{0,0395}
\setval{mtv-f6a1-frec-sd}{0,0011}
\setval{mtv-m6-frec-mean}{0,0565}
\setval{mtv-m6-frec-sd}{0,0003}
\setval{mtv-6-frec-diff}{\ensuremath{-0,0170}}  % f6a1 minus m6
\setval{mtv-6-frec-rel}{\ensuremath{-30,2}}  % relative difference, %
\setval{mtv-6-frec-p}{\ensuremath{<}0,001}  % Welch, two-sided
\setval{mtv-6-frec-pt}{\ensuremath{p<0,001}}
\setval{mtv-6-frec-diffabs}{0,0170}
\setval{mtv-6-frec-relabs}{30,2}
\setval{mtv-6-frec-sdx}{15}  % difference in units of the larger seed SD
\setval{twin-maxdiff}{0,0002}  % largest change of test Hit@10 when the twin edges are left out, any run
\setval{main-hit-sd-max}{0,004}  % largest seed SD of test Hit@10 among the main runs

% --- GAT and popularity --------------------------------------------------------------------------
\setval{gat-hit-mean}{0,787}
\setval{gat-hit-sd}{0,004}
\setval{pop-hit}{0,423}
\setval{gat-ndcg-mean}{0,553}
\setval{gat-ndcg-sd}{0,005}
\setval{pop-ndcg}{0,249}
\setval{gat-fhit-mean}{0,0476}
\setval{gat-fhit-sd}{0,0049}
\setval{pop-fhit}{0,0094}
\setval{gat-fndcg-mean}{0,0246}
\setval{gat-fndcg-sd}{0,0027}
\setval{pop-fndcg}{0,0049}
\setval{gat-frec-mean}{0,0897}
\setval{gat-frec-sd}{0,0082}
\setval{pop-frec}{0,0172}
\setval{gat-val-mean}{0,1780}  % validation BPR, mean of the 3 seeds
\setval{gat-val-sd}{0,0031}

% --- By the user's number of training edges: Hit@10, mean over the seeds run (HAN logs) ----------
\setval{testdeg-d0}{8}  % test edges in the group
\setval{testdeg-d1}{15.831}  % test edges in the group
\setval{testdeg-d2}{27.948}  % test edges in the group
\setval{testdeg-d3}{38.197}  % test edges in the group
\setval{han-m1-hit-d1}{0,571}
\setval{han-m1-hit-d2}{0,629}
\setval{han-m1-hit-d3}{0,666}
\setval{han-m2-hit-d1}{0,577}
\setval{han-m2-hit-d2}{0,646}
\setval{han-m2-hit-d3}{0,683}
\setval{han-m3-hit-d1}{0,576}
\setval{han-m3-hit-d2}{0,640}
\setval{han-m3-hit-d3}{0,683}
\setval{han-m4-hit-d1}{0,572}
\setval{han-m4-hit-d2}{0,619}
\setval{han-m4-hit-d3}{0,639}
\setval{han-m5-hit-d1}{0,596}
\setval{han-m5-hit-d2}{0,662}
\setval{han-m5-hit-d3}{0,693}
\setval{han-m6-hit-d1}{0,622}
\setval{han-m6-hit-d2}{0,686}
\setval{han-m6-hit-d3}{0,715}
\setval{han-m7-hit-d1}{0,605}
\setval{han-m7-hit-d2}{0,668}
\setval{han-m7-hit-d3}{0,698}
\setval{han-m8-hit-d1}{0,597}
\setval{han-m8-hit-d2}{0,658}
\setval{han-m8-hit-d3}{0,694}
\setval{han-r4d1-hit-d1}{0,601}
\setval{han-r4d1-hit-d2}{0,666}
\setval{han-r4d1-hit-d3}{0,705}
\setval{han-r4d2-hit-d1}{0,595}
\setval{han-r4d2-hit-d2}{0,658}
\setval{han-r4d2-hit-d3}{0,693}
\setval{han-r4d3-hit-d1}{0,609}
\setval{han-r4d3-hit-d2}{0,673}
\setval{han-r4d3-hit-d3}{0,709}
\setval{han-r4d4-hit-d1}{0,604}
\setval{han-r4d4-hit-d2}{0,668}
\setval{han-r4d4-hit-d3}{0,701}
\setval{han-r4d5-hit-d1}{0,598}
\setval{han-r4d5-hit-d2}{0,655}
\setval{han-r4d5-hit-d3}{0,689}
\setval{han-r6d1-hit-d1}{0,605}
\setval{han-r6d1-hit-d2}{0,670}
\setval{han-r6d1-hit-d3}{0,703}
\setval{han-r6d2-hit-d1}{0,639}
\setval{han-r6d2-hit-d2}{0,700}
\setval{han-r6d2-hit-d3}{0,739}
\setval{han-r6d3-hit-d1}{0,629}
\setval{han-r6d3-hit-d2}{0,692}
\setval{han-r6d3-hit-d3}{0,724}
\setval{han-r6d4-hit-d1}{0,610}
\setval{han-r6d4-hit-d2}{0,676}
\setval{han-r6d4-hit-d3}{0,702}
\setval{han-r6d5-hit-d1}{0,640}
\setval{han-r6d5-hit-d2}{0,702}
\setval{han-r6d5-hit-d3}{0,736}
\setval{han-f4a1-hit-d1}{0,562}
\setval{han-f4a1-hit-d2}{0,631}
\setval{han-f4a1-hit-d3}{0,668}
\setval{han-f4a2-hit-d1}{0,579}
\setval{han-f4a2-hit-d2}{0,645}
\setval{han-f4a2-hit-d3}{0,683}
\setval{han-f4c-hit-d1}{0,584}
\setval{han-f4c-hit-d2}{0,653}
\setval{han-f4c-hit-d3}{0,694}
\setval{han-f6a1-hit-d1}{0,592}
\setval{han-f6a1-hit-d2}{0,663}
\setval{han-f6a1-hit-d3}{0,700}
\setval{han-f6a2-hit-d1}{0,596}
\setval{han-f6a2-hit-d2}{0,661}
\setval{han-f6a2-hit-d3}{0,699}
\setval{han-f6c-hit-d1}{0,595}
\setval{han-f6c-hit-d2}{0,662}
\setval{han-f6c-hit-d3}{0,698}
\setval{mtv-m1-hit-d1}{0,514}
\setval{mtv-m1-hit-d2}{0,591}
\setval{mtv-m1-hit-d3}{0,641}
\setval{mtv-m2-hit-d1}{0,581}
\setval{mtv-m2-hit-d2}{0,644}
\setval{mtv-m2-hit-d3}{0,657}
\setval{mtv-m3-hit-d1}{0,570}
\setval{mtv-m3-hit-d2}{0,634}
\setval{mtv-m3-hit-d3}{0,639}
\setval{mtv-m4-hit-d1}{0,097}
\setval{mtv-m4-hit-d2}{0,096}
\setval{mtv-m4-hit-d3}{0,093}
\setval{mtv-m5-hit-d1}{0,558}
\setval{mtv-m5-hit-d2}{0,643}
\setval{mtv-m5-hit-d3}{0,671}
\setval{mtv-m6-hit-d1}{0,582}
\setval{mtv-m6-hit-d2}{0,647}
\setval{mtv-m6-hit-d3}{0,654}
\setval{mtv-m7-hit-d1}{0,552}
\setval{mtv-m7-hit-d2}{0,644}
\setval{mtv-m7-hit-d3}{0,659}
\setval{mtv-m8-hit-d1}{0,561}
\setval{mtv-m8-hit-d2}{0,618}
\setval{mtv-m8-hit-d3}{0,643}
\setval{mtv-r4d1-hit-d1}{0,389}
\setval{mtv-r4d1-hit-d2}{0,562}
\setval{mtv-r4d1-hit-d3}{0,615}
\setval{mtv-r4d2-hit-d1}{0,439}
\setval{mtv-r4d2-hit-d2}{0,589}
\setval{mtv-r4d2-hit-d3}{0,638}
\setval{mtv-r4d3-hit-d1}{0,541}
\setval{mtv-r4d3-hit-d2}{0,622}
\setval{mtv-r4d3-hit-d3}{0,624}
\setval{mtv-r4d4-hit-d1}{0,552}
\setval{mtv-r4d4-hit-d2}{0,635}
\setval{mtv-r4d4-hit-d3}{0,662}
\setval{mtv-r4d5-hit-d1}{0,358}
\setval{mtv-r4d5-hit-d2}{0,539}
\setval{mtv-r4d5-hit-d3}{0,633}
\setval{mtv-r6d1-hit-d1}{0,490}
\setval{mtv-r6d1-hit-d2}{0,578}
\setval{mtv-r6d1-hit-d3}{0,625}
\setval{mtv-r6d2-hit-d1}{0,500}
\setval{mtv-r6d2-hit-d2}{0,635}
\setval{mtv-r6d2-hit-d3}{0,663}
\setval{mtv-r6d3-hit-d1}{0,546}
\setval{mtv-r6d3-hit-d2}{0,629}
\setval{mtv-r6d3-hit-d3}{0,650}
\setval{mtv-r6d4-hit-d1}{0,561}
\setval{mtv-r6d4-hit-d2}{0,647}
\setval{mtv-r6d4-hit-d3}{0,650}
\setval{mtv-r6d5-hit-d1}{0,480}
\setval{mtv-r6d5-hit-d2}{0,615}
\setval{mtv-r6d5-hit-d3}{0,654}
\setval{mtv-f4a1-hit-d1}{0,485}
\setval{mtv-f4a1-hit-d2}{0,599}
\setval{mtv-f4a1-hit-d3}{0,648}
\setval{mtv-f4a2-hit-d1}{0,533}
\setval{mtv-f4a2-hit-d2}{0,613}
\setval{mtv-f4a2-hit-d3}{0,635}
\setval{mtv-f4c-hit-d1}{0,535}
\setval{mtv-f4c-hit-d2}{0,629}
\setval{mtv-f4c-hit-d3}{0,662}
\setval{mtv-f6a1-hit-d1}{0,446}
\setval{mtv-f6a1-hit-d2}{0,590}
\setval{mtv-f6a1-hit-d3}{0,648}
\setval{mtv-f6a2-hit-d1}{0,456}
\setval{mtv-f6a2-hit-d2}{0,591}
\setval{mtv-f6a2-hit-d3}{0,649}
\setval{mtv-f6c-hit-d1}{0,466}
\setval{mtv-f6c-hit-d2}{0,599}
\setval{mtv-f6c-hit-d3}{0,653}
\setval{gat-hit-d1}{0,755}
\setval{gat-hit-d2}{0,798}
\setval{gat-hit-d3}{0,793}
\setval{pop-hit-d1}{0,431}
\setval{pop-hit-d2}{0,435}
\setval{pop-hit-d3}{0,412}
\setval{han-4-hit-d1-diff}{\ensuremath{-0,034}}
\setval{han-4-hit-d2-diff}{\ensuremath{-0,031}}
\setval{han-4-hit-d3-diff}{\ensuremath{-0,025}}
\setval{han-6-hit-d1-diff}{\ensuremath{-0,031}}
\setval{han-6-hit-d2-diff}{\ensuremath{-0,023}}
\setval{han-6-hit-d3-diff}{\ensuremath{-0,015}}
\setval{gat-lead-d1}{\ensuremath{+0,117}}  % GAT minus the HAN validation pick
\setval{gat-lead-d2}{\ensuremath{+0,098}}  % GAT minus the HAN validation pick
\setval{gat-lead-d3}{\ensuremath{+0,054}}  % GAT minus the HAN validation pick

% --- All 24 sets at seed 1: families, order, agreement (HAN notes, Part 10) ----------------------
\setval{han-r4-hit-mean}{0,668}
\setval{han-r4-hit-min}{0,660}
\setval{han-r4-hit-max}{0,677}
\setval{han-r4-fhit-mean}{0,0234}
\setval{han-r4-fhit-min}{0,0203}
\setval{han-r4-fhit-max}{0,0267}
\setval{han-4-pos-hit}{thứ 6, 7 và 8}  % places of the searched sets among the 8 candidate-drawn sets
\setval{han-4-pos-val}{thứ 6, 7 và 8}  % places of the searched sets among the 8 candidate-drawn sets
\setval{han-4-pos-fhit}{thứ 3, 4 và 7}  % places of the searched sets among the 8 candidate-drawn sets
\setval{han-r6-hit-mean}{0,691}
\setval{han-r6-hit-min}{0,673}
\setval{han-r6-hit-max}{0,706}
\setval{han-r6-fhit-mean}{0,0264}
\setval{han-r6-fhit-min}{0,0228}
\setval{han-r6-fhit-max}{0,0282}
\setval{han-6-pos-hit}{thứ 6, 7 và 8}  % places of the searched sets among the 8 candidate-drawn sets
\setval{han-6-pos-val}{thứ 6, 7 và 8}  % places of the searched sets among the 8 candidate-drawn sets
\setval{han-6-pos-fhit}{thứ 5, 6 và 7}  % places of the searched sets among the 8 candidate-drawn sets
\setval{mtv-r4-hit-mean}{0,585}
\setval{mtv-r4-hit-min}{0,548}
\setval{mtv-r4-hit-max}{0,631}
\setval{mtv-r4-fhit-mean}{0,0203}
\setval{mtv-r4-fhit-min}{0,0107}
\setval{mtv-r4-fhit-max}{0,0316}
\setval{mtv-4-pos-hit}{thứ 2, 3 và 5}  % places of the searched sets among the 8 candidate-drawn sets
\setval{mtv-4-pos-val}{thứ 2, 4 và 5}  % places of the searched sets among the 8 candidate-drawn sets
\setval{mtv-4-pos-fhit}{thứ 1, 3 và 4}  % places of the searched sets among the 8 candidate-drawn sets
\setval{mtv-r6-hit-mean}{0,613}
\setval{mtv-r6-hit-min}{0,583}
\setval{mtv-r6-hit-max}{0,632}
\setval{mtv-r6-fhit-mean}{0,0227}
\setval{mtv-r6-fhit-min}{0,0128}
\setval{mtv-r6-fhit-max}{0,0391}
\setval{mtv-6-pos-hit}{thứ 5, 6 và 7}  % places of the searched sets among the 8 candidate-drawn sets
\setval{mtv-6-pos-val}{thứ 4, 5 và 7}  % places of the searched sets among the 8 candidate-drawn sets
\setval{mtv-6-pos-fhit}{thứ 2, 3 và 4}  % places of the searched sets among the 8 candidate-drawn sets
\setval{han-r4-vs-manual-diff}{\ensuremath{+0,005}}  % mean of the 5 draws (seed 1) minus the manual set (3 seeds)
\setval{han-r4-vs-manual-rel}{\ensuremath{+0,7}}  % %
\setval{han-r6-vs-manual-diff}{\ensuremath{+0,004}}  % mean of the 5 draws (seed 1) minus the manual set (3 seeds)
\setval{han-r6-vs-manual-rel}{\ensuremath{+0,6}}  % %
\setval{han-valtest-rho}{\ensuremath{-0,94}}
\setval{han-valfull-rho}{\ensuremath{-0,71}}
\setval{han-mtv-rho}{\ensuremath{+0,09}}  % HAN against metapath2vec over the 23 sets (M4 left out)
\setval{han-wins-sampled}{24}  % of 24 sets
\setval{mtv-wins-full}{14}  % of 24 sets
\setval{han-valpick}{r6d2}  % lowest validation BPR among the 24 HAN sets
\setval{mtv-valpick}{m5}  % lowest validation BPR among the 23 metapath2vec sets
\setval{han-4-dense-min}{0,660}
\setval{han-4-dense-max}{0,677}
\setval{han-4-nodense-min}{0,635}
\setval{han-4-nodense-max}{0,658}
\setval{han-6-dense-min}{0,673}
\setval{han-6-dense-max}{0,706}
\setval{han-6-nodense-min}{0,666}
\setval{han-6-nodense-max}{0,666}
\setval{han-4-span}{0,042}  % range of HAN test Hit@10 over all 4-path sets, seed 1
\setval{han-4-nsets}{11}
\setval{han-searched4-mean-s1}{0,648}  % mean of the three searched sets, seed 1
\setval{han-6-span}{0,041}  % range of HAN test Hit@10 over all 6-path sets, seed 1
\setval{han-6-nsets}{9}
\setval{han-searched6-mean-s1}{0,666}  % mean of the three searched sets, seed 1
\setval{searched-with-pupup}{6}  % searched sets holding PUPUP
\setval{searched-with-puppup}{4}  % searched sets holding PUPPUP
\setval{searched-empty-min}{1}  % near-empty category paths per searched set
\setval{searched-empty-max}{3}

% --- What HAN does with the paths: weights at seed 1 against the search (LMSPS part 4b) ----------
\setval{kind-users-paths}{12}
\setval{kind-users-rank}{12,7}
\setval{kind-users-uses}{46}
\setval{kind-users-dropped}{0}
\setval{kind-users-weight}{1,34}
\setval{kind-empty-paths}{6}
\setval{kind-empty-rank}{11,8}
\setval{kind-empty-uses}{25}
\setval{kind-empty-dropped}{16}
\setval{kind-empty-weight}{0,00}
\setval{kind-cat-paths}{10}
\setval{kind-cat-rank}{23,8}
\setval{kind-cat-uses}{22}
\setval{kind-cat-dropped}{3}
\setval{kind-cat-weight}{1,12}
\setval{kind-pp-paths}{4}
\setval{kind-pp-rank}{28,8}
\setval{kind-pp-uses}{16}
\setval{kind-pp-dropped}{2}
\setval{kind-pp-weight}{0,85}
\setval{han-dropped}{21}  % weights below 0.01, 24 sets
\setval{weight-drop}{0,01}  % a path weight below this counts as dropped
\setval{han-empty-kept-min}{0,26}  % largest near-empty-path weight (beta) in a searched set, f4c aside
\setval{han-empty-kept-max}{0,31}
\setval{han-empty-kept-f4c}{0,06}
\setval{weight-pupup}{1,76}  % mean weight of PUPUP as a multiple of an equal share
\setval{uses-pupup}{10}
\setval{weight-puppup}{1,31}  % mean weight of PUPPUP as a multiple of an equal share
\setval{uses-puppup}{5}
\setval{weight-pup}{1,45}  % mean weight of PUP as a multiple of an equal share
\setval{uses-pup}{8}
\setval{weight-pcp}{1,04}  % mean weight of PCP as a multiple of an equal share
\setval{uses-pcp}{7}
\setval{weight-ppccpp}{0,00}  % mean weight of PPCCPP as a multiple of an equal share
\setval{uses-ppccpp}{6}
\setval{pref-cand-rho}{\ensuremath{+0,54}}  % search preference against han, n = 16
\setval{pref-cand-p}{0,035}  % permutation p, 20,000 shuffles, Random(0)
\setval{pref-cand-pt}{\ensuremath{p=0,035}}
\setval{pref-cand-n}{16}
\setval{pref-cand-val-rho}{\ensuremath{-0,61}}  % search preference against val, n = 16
\setval{pref-cand-val-p}{0,017}  % permutation p, 20,000 shuffles, Random(0)
\setval{pref-cand-val-pt}{\ensuremath{p=0,017}}
\setval{pref-cand-val-n}{16}
\setval{pref-cand-full-rho}{\ensuremath{+0,31}}  % search preference against full, n = 16
\setval{pref-cand-full-p}{0,24}  % permutation p, 20,000 shuffles, Random(0)
\setval{pref-cand-full-pt}{\ensuremath{p=0,24}}
\setval{pref-cand-full-n}{16}
\setval{pref-random-rho}{\ensuremath{+0,28}}  % search preference against han, n = 10
\setval{pref-random-p}{0,43}  % permutation p, 20,000 shuffles, Random(0)
\setval{pref-random-pt}{\ensuremath{p=0,43}}
\setval{pref-random-n}{10}
\setval{pref-all-rho}{\ensuremath{+0,18}}  % search preference against han, n = 24
\setval{pref-all-p}{0,39}  % permutation p, 20,000 shuffles, Random(0)
\setval{pref-all-pt}{\ensuremath{p=0,39}}
\setval{pref-all-n}{24}
\setval{pref-mtv-rho}{\ensuremath{+0,29}}  % search preference against m2v, n = 23
\setval{pref-mtv-p}{0,19}  % permutation p, 20,000 shuffles, Random(0)
\setval{pref-mtv-pt}{\ensuremath{p=0,19}}
\setval{pref-mtv-n}{23}

% --- Where the runs differ: Hit@10 by product group, seed 1 (HAN review/evidence/by_group.txt) ---
\setval{bg-pccp-reach-products-pct}{18,9}  % products a P-C-C-P walk reaches
\setval{bg-pccp-reach-test-pct}{27,5}  % test edges whose product it reaches
\setval{bg-4-gap-no}{\ensuremath{-0,023}}  % searched sets minus the sets with a dense category path
\setval{bg-4-gapabs-no}{0,023}
\setval{bg-4-gap-reach}{\ensuremath{-0,010}}  % searched sets minus the sets with a dense category path
\setval{bg-4-gapabs-reach}{0,010}
\setval{bg-4-gap-book}{\ensuremath{-0,026}}  % searched sets minus the sets with a dense category path
\setval{bg-4-gapabs-book}{0,026}
\setval{bg-4-gap-music}{\ensuremath{-0,017}}  % searched sets minus the sets with a dense category path
\setval{bg-4-gapabs-music}{0,017}
\setval{bg-4-gap-dvd}{\ensuremath{-0,009}}  % searched sets minus the sets with a dense category path
\setval{bg-4-gapabs-dvd}{0,009}
\setval{bg-4-gap-video}{\ensuremath{-0,014}}  % searched sets minus the sets with a dense category path
\setval{bg-4-gapabs-video}{0,014}
\setval{bg-4-a1gap-no}{\ensuremath{-0,034}}  % a1 minus the manual set
\setval{bg-4-a1gap-reach}{\ensuremath{-0,019}}  % a1 minus the manual set
\setval{bg-6-gap-no}{\ensuremath{-0,034}}  % searched sets minus the sets with a dense category path
\setval{bg-6-gapabs-no}{0,034}
\setval{bg-6-gap-reach}{\ensuremath{-0,002}}  % searched sets minus the sets with a dense category path
\setval{bg-6-gapabs-reach}{0,002}
\setval{bg-6-gap-book}{\ensuremath{-0,030}}  % searched sets minus the sets with a dense category path
\setval{bg-6-gapabs-book}{0,030}
\setval{bg-6-gap-music}{\ensuremath{-0,034}}  % searched sets minus the sets with a dense category path
\setval{bg-6-gapabs-music}{0,034}
\setval{bg-6-gap-dvd}{\ensuremath{-0,004}}  % searched sets minus the sets with a dense category path
\setval{bg-6-gapabs-dvd}{0,004}
\setval{bg-6-gap-video}{\ensuremath{-0,007}}  % searched sets minus the sets with a dense category path
\setval{bg-6-gapabs-video}{0,007}
\setval{bg-6-a1gap-no}{\ensuremath{-0,031}}  % a1 minus the manual set
\setval{bg-6-a1gap-reach}{\ensuremath{-0,011}}  % a1 minus the manual set
\setval{bg-users}{5.000}  % test users whose full-ranking top 10 is examined
\setval{bg-all-deg}{13,4}  % mean number of training users of a product
\setval{bg-han-deg-med}{74}  % median over the sets, M4 aside
\setval{bg-han-deg-min}{57}
\setval{bg-han-deg-max}{82}
\setval{bg-mtv-deg-med}{6}
\setval{bg-gat-deg}{152}
\setval{bg-pop-deg}{550}
\setval{bg-han-top1-med}{24}  % % of the top 10 from the 1% most-bought
\setval{bg-mtv-top1-max}{0,1}
\setval{bg-gat-top1}{51}
\setval{bg-han-mtv-book}{\ensuremath{+0,015}}  % HAN minus metapath2vec, median over the 23 sets
\setval{bg-han-mtv-music}{\ensuremath{-0,047}}  % HAN minus metapath2vec, median over the 23 sets
\setval{bg-han-mtv-dvd}{\ensuremath{+0,282}}  % HAN minus metapath2vec, median over the 23 sets
\setval{bg-han-mtv-video}{\ensuremath{+0,216}}  % HAN minus metapath2vec, median over the 23 sets

% --- Training and cost (HAN logs; HAN notes, Parts 9-10) -----------------------------------------
\setval{han-runs}{32}
\setval{han-run-min}{9,0}  % median minutes per run
\setval{han-hours}{5,0}  % GPU hours, all runs
\setval{mtv-runs}{32}
\setval{mtv-run-min}{5,2}  % median minutes per run
\setval{mtv-hours}{2,9}  % GPU hours, all runs
\setval{gat-runs}{3}
\setval{gat-run-min}{12,8}  % median minutes per run
\setval{gat-hours}{0,6}  % GPU hours, all runs
\setval{models-hours}{9}  % HAN, metapath2vec and GAT
\setval{han-stop-min}{63}  % epoch (metapath2vec: pass) of the early stop
\setval{han-stop-max}{112}
\setval{gat-stop-min}{130}  % epoch (metapath2vec: pass) of the early stop
\setval{gat-stop-max}{157}
\setval{mtv-stop-min}{1,5}  % epoch (metapath2vec: pass) of the early stop
\setval{mtv-stop-max}{3,0}
\setval{gat-cap}{150}  % epoch cap of GAT
\setval{gat-rerun-seed}{2}  % the GAT seed that hit the cap and was rerun
\setval{gat-rerun-cap}{300}
\setval{gat-rerun-stop}{157}
\setval{han-searched6-min-lo}{11,5}  % minutes per run, the three searched 6-path sets, seed 1
\setval{han-searched6-min-hi}{11,9}
\setval{han-other6-min-lo}{10,1}  % the other 6-path sets (M6, random), seed 1
\setval{han-other6-min-hi}{13,7}
\setval{mtv-searched6-min-lo}{3,4}  % minutes per run, the three searched 6-path sets, seed 1
\setval{mtv-searched6-min-hi}{3,4}
\setval{mtv-other6-min-lo}{3,3}  % the other 6-path sets (M6, random), seed 1
\setval{mtv-other6-min-hi}{10,0}

% --- Repeatability: a same-seed repeat of one HAN run on the same GPU (logs/addition runs) -------
\setval{repeat-han-best}{88}  % best epoch, the same in both runs (HAN, a1 with 4 paths, seed 1)
\setval{repeat-han-val}{0,3056}  % validation BPR, the same in both runs
\setval{repeat-han-hit}{0,6348}  % test Hit@10, the same in both runs
\setval{repeat-han-groupdiff}{0,0001}  % largest Hit@10 difference in a user-degree group
\setval{han-params-min}{27,8}  % millions
\setval{han-params-max}{28,3}  % millions
\setval{han-user-params}{27,4}  % millions: one 512-d vector per user
\setval{han-gain-med}{0,0006}  % validation BPR gained in the last 10 epochs
\setval{han-gain-max}{0,0037}
\setval{gat-gain-min}{0,0007}
\setval{gat-gain-max}{0,0022}
\setval{searched-gain-max}{0,0006}

% --- Side check: HAN with a 5-epoch warm-up, seed 1 (HAN notes, Part 10) -------------------------
\setval{wu5-warmup}{5}  % epochs of path-weight warm-up in the side check (asserted in every log)
\setval{wu5-diff-mean}{\ensuremath{-0,017}}
\setval{wu5-diff-min}{\ensuremath{-0,066}}
\setval{wu5-diff-max}{\ensuremath{+0,007}}
\setval{wu5-dropped}{41}
\setval{wu5-rho}{0,64}
\setval{wu5-4-pos}{thứ 4, 5 và 7}
\setval{wu5-6-pos}{thứ 4, 5 và 7}
\setval{wu5-f6a1-hit}{0,663}
\setval{wu5-m6-hit}{0,644}
\setval{wu5-m6-kept}{3}  % of its 6 paths

% --- From other sources (typed in make_numbers.py, with the source) ------------------------------
\setval{mcauley-raw-reviews}{142,8 triệu}  % McAuley's Amazon product data page, 2014 version (cseweb.ucsd.edu/~jmcauley/datasets/amazon/links.html), read 5 November 2026: raw reviews
\setval{mcauley-dedup-reviews}{83,68 triệu}  % the same page: reviews left after removing duplicate items
\setval{lmsps-t4-gain-min}{1,3}  % LMSPS paper, Table 4: Micro-F1 of HAN* minus HAN, four datasets
\setval{lmsps-t4-gain-max}{2,4}
\setval{dacrema-algos}{18}  % Dacrema et al., RecSys 2019, abstract: neural recommenders analysed
\setval{dacrema-repro}{7}  % the same: reproduced with reasonable effort
\setval{dacrema-beaten}{6}  % the same: often outperformed by simple heuristics
\setval{diffmg-gems-min}{800}  % DiffMG paper, Table 5: GEMS search cost, GPU minutes
\setval{diffmg-gems-max}{2.000}
\setval{diffmg-own-max}{1,6}  % the same table: DiffMG search cost (three runs)

% --- Provenance: the checkouts these numbers were computed from ----------------------------------
\setval{repo-lmsps-commit}{\detokenize{5e32eb6}}  % read from .git/HEAD
\setval{repo-lmsps-branch}{\detokenize{claude/fix-amazon-extraction}}
\setval{repo-han-commit}{\detokenize{2ea02f6}}  % read from .git/HEAD
\setval{repo-han-branch}{\detokenize{claude/han-review}}


% The paper's numbers. \ref*{P-...} prints them without a link, as the targets are in the other
% PDF. \addFileDependency (Overleaf's recipe for xr) tells latexmk that this document needs
% paper_v2_vn.aux, so the rule in latexmkrc builds the paper first when it is missing or old.
\usepackage{xr}
\makeatletter
\newcommand*{\addFileDependency}[1]{%
  \typeout{(#1)}%
  \@addtofilelist{#1}%
  \IfFileExists{#1}{}{\typeout{No file #1.}}}
\makeatother
\externaldocument[P-]{paper_v2_vn}
\addFileDependency{paper_v2_vn.tex}
\addFileDependency{paper_v2_vn.aux}

% This document's tables and figure: Bảng S1, Hình S1, ...
\renewcommand{\thetable}{S\arabic{table}}
\renewcommand{\thefigure}{S\arabic{figure}}

\title[Tài liệu bổ sung]{Tài liệu bổ sung}
\subtitle{cho bài báo ``Đánh giá chi phí và hiệu quả của việc tìm kiếm
  meta-path\texorpdfstring{\\}{ } cho gợi ý sản phẩm dựa trên HAN''}
\author{Phan Hoàng Tuấn Anh}
\email{phananh85327@gmail.com}
\affiliation{%
  \department{Khoa Công nghệ Thông tin}
  \institution{Trường Đại học Công Thương Thành phố Hồ Chí Minh (HUIT)}
  \city{Thành phố Hồ Chí Minh}
  \country{Việt Nam}}

\begin{document}

\maketitle
\hypersetup{pdftitle={Tài liệu bổ sung: Đánh giá chi phí và hiệu quả của việc tìm kiếm meta-path
  cho gợi ý sản phẩm dựa trên HAN}, pdfauthor={Phan Hoàng Tuấn Anh}}

Tài liệu này chứa những chi tiết mà bài báo không trình bày. Các mục của tài liệu được đánh số bằng
chữ cái (A--H), còn các bảng và hình mang tiền tố S. Những số mục, bảng và hình không có tiền tố
(chẳng hạn Mục~\ref*{P-sec:data} hay Bảng~\ref*{P-tab:main}) chỉ đến bài báo.

\appendix
% sections/appendix.tex -- the body of the supplementary document, supplement_vn.tex (paper_v2_vn,
% translated in L8 from paper_v2's sections/appendix.tex; terms as in glossary_vn.tex). Until N31
% it was the paper's appendix. Its references to the paper are \ref*{P-...} (xr, see
% supplement_vn.tex), and its own sections are called Mục, lettered A-H. Its sections follow the
% paper: the study at a glance, the data (Section 4), changes from the base papers and the search
% (Section 5), the random draws and the settings (Section 6), full results (Section 7),
% reproducibility. One addition to paper_v2: the Records paragraph names make_numbers.py's
% --lang vi option, which writes this version's numbers. Every number is \val{} from numbers.tex.
\section{Tổng quan nghiên cứu}
\label{app:overview}

Hình~\ref{fig:pipeline} trình bày nghiên cứu từ dữ liệu thô đến các cạnh thử nghiệm
(Mục~\ref*{P-sec:data}--\ref*{P-sec:setup}).

% f1_pipeline.tex -- F1: the study at a glance (paper_v2_vn, translated in L3 from paper_v2's
% figures/f1_pipeline.tex; TikZ, no data file). Colours as in paper_v2_vn.tex (searched cSearch,
% manual cManual, random cRandom).
\begin{widefigure}[htbp]
  \centering
  \resizebox{\linewidth}{!}{%
  \begin{tikzpicture}[
      font=\footnotesize,
      box/.style={draw=black!45, rounded corners=2pt, align=center, inner sep=4pt,
        text width=#1, minimum height=11mm, execute at begin node={\hyphenpenalty=10000}},
      box/.default=36mm,
      way/.style={box=44mm, draw=#1, fill=#1!10, line width=0.7pt},
      flow/.style={-{Stealth[length=2.2mm]}, black!55, line width=0.6pt}]
    % data
    \node[box] (raw) {Amazon-Meta (SNAP)\\\val{raw-records} sản phẩm,\\\val{raw-reviews} dòng nhận
      xét};
    \node[box, below=4mm of raw] (clean) {Làm sạch: gộp ấn bản, loại \val{placeholder-ids} tài
      khoản, giữ nhận xét tích cực, \val{rule-kcore}-lõi};
    \node[box, below=4mm of clean] (graph) {Đồ thị: sản phẩm, danh mục, người dùng. Cạnh người
      dùng--sản phẩm chia 80/10/10};
    % three ways to choose the paths
    \node[way=cManual, right=12mm of raw] (hand) {\textbf{Thủ công}: các tập M1--M8, thiết kế với
      sự hỗ trợ của AI, cố định từ trước};
    \node[way=cSearch, below=4mm of hand] (search) {\textbf{Tìm kiếm}: LMSPS cho dự đoán liên kết,
      \val{search-runs} lần chạy cho mỗi kích thước tập, kết hợp thành \code{a1}, \code{a2},
      \code{c}};
    \node[way=cRandom, below=4mm of search] (random) {\textbf{Ngẫu nhiên (đối chứng)}: 5 lần chọn cho mỗi kích
      thước từ cùng \val{search-cands} ứng viên};
    % models and test
    \node[box=36mm, right=12mm of search] (models) {HAN và metapath2vec, huấn luyện trên từng tập
      trong 24 tập meta-path};
    \node[box=36mm, right=10mm of models] (test) {Cạnh thử nghiệm: Hit@10 lấy mẫu (cố định từ
      trước), NDCG@10, xếp hạng toàn bộ, nhóm người dùng};
    \node[box=36mm, below=4mm of models, yshift=-6mm] (refs) {GAT trên toàn bộ đồ thị (bỏ qua
      kiểu) và mô hình phổ biến};
    % arrows
    \draw[flow] (raw) -- (clean);
    \draw[flow] (clean) -- (graph);
    \draw[flow] (graph.east) -- ++(4mm,0) |- (hand.west);
    \draw[flow] (graph.east) -- ++(4mm,0) |- (search.west);
    \draw[flow] (graph.east) -- ++(4mm,0) |- (random.west);
    \draw[flow] (hand.east) -- ++(4mm,0) |- (models.west);
    \draw[flow] (search.east) -- (models.west);
    \draw[flow] (random.east) -- ++(4mm,0) |- (models.west);
    \draw[flow] (models) -- (test);
    \draw[flow] (graph.south) -- ++(0,-6mm) -| ([xshift=-6mm]refs.west) -- (refs.west);
    \draw[flow] (refs.east) -| (test.south);
  \end{tikzpicture}}
  \Description{Sơ đồ luồng. Dữ liệu Amazon-Meta được làm sạch và chuyển thành đồ thị. Hai hộp, thủ công và tìm kiếm, chọn các tập meta-path từ cùng các ứng viên, cùng một hộp ngẫu nhiên làm đối chứng. Mọi tập được chuyển cho HAN và metapath2vec. GAT và mô hình phổ biến dùng toàn bộ đồ thị. Mọi mô hình được thử nghiệm trên cùng các cạnh giữ lại.}
  \caption{Tổng quan nghiên cứu. Từ cùng các meta-path ứng viên, các tập meta-path cho HAN được chọn
    thủ công hoặc bằng tìm kiếm, cùng các tập chọn ngẫu nhiên làm đối chứng. Mọi tập được huấn luyện
    trên HAN và metapath2vec rồi thử nghiệm trên cùng các cạnh giữ lại, với GAT và mô hình phổ biến
    làm mô hình cơ sở không dùng meta-path.}
  \label{fig:pipeline}
\end{widefigure}


\section{Dữ liệu và đồ thị}
\label{app:graph}

Bảng~\ref{tab:graph} tóm tắt đồ thị và phép chia ở Mục~\ref*{P-sec:data}. Phần còn lại của mục
này nêu những chi tiết về dữ liệu mà Mục~\ref*{P-sec:data} không trình bày.

\paragraph{Gộp ấn bản.} Quy tắc gộp nối \val{edition-pairs} cặp thành \val{edition-clusters} cụm
gồm \val{edition-size-min} đến \val{edition-size-max} sản phẩm (\val{edition-size-two-pct}\% số cụm
chỉ gồm hai sản phẩm) và để lại \val{raw-works} tác phẩm từ \val{raw-usable} bản ghi sản phẩm dùng
được. Với dữ liệu nhận xét Amazon ra đời sau và lớn hơn của \citet{mcauley2015image}, trang dữ liệu
cảnh báo rằng các nhận xét thô có chứa bản trùng lặp. Các tệp đã loại trùng lặp giữ
\val{mcauley-dedup-reviews} trong số \val{mcauley-raw-reviews}
nhận xét.\footnote{\url{https://cseweb.ucsd.edu/~jmcauley/datasets/amazon/links.html}, truy cập
ngày 5 tháng 11 năm 2026.}

\paragraph{Sản phẩm song sinh.} Một phép kiểm tra thứ hai tìm các sản phẩm có tập người nhận xét
giống hệt nhau. Phép kiểm tra này tìm thấy \val{twin-products} sản phẩm trong \val{twin-groups}
nhóm, với \val{twin-test} cạnh thử nghiệm có một sản phẩm song sinh như vậy trong tập huấn luyện.
Loại bỏ các cạnh này không làm Hit@10 thử nghiệm của lần chạy nào thay đổi quá \val{twin-maxdiff}.

\paragraph{Đặc trưng tiêu đề.} Vectơ tiêu đề của một sản phẩm là một vectơ LSA: TF-IDF trên bộ từ
vựng gồm \val{title-vocab} token, giảm chiều bằng SVD rút gọn xuống \val{title-dim} chiều. Vectơ này
khác 0 với \val{title-coverage-pct}\% số sản phẩm.

\paragraph{Liên kết danh mục.} Sản phẩm chỉ liên kết với các danh mục lá của nó. Nếu liên kết thêm
mọi danh mục cấp trên, sản phẩm trung vị sẽ có \val{cat-all-median} danh mục thay vì
\val{after-leafcats-med}, \val{cat-all-largest-pct}\% số sản phẩm sẽ rơi vào cùng một danh mục và
các đường đi P--C--P sẽ dày đặc gấp \val{cat-pcp-ratio} lần. Với liên kết đến danh mục lá,
\code{PCP} chỉ để \val{pcp-isolated} sản phẩm không có láng giềng nào.

% t2_graph.tex -- T2: the graph and the split (paper_v2_vn, translated in L3 from paper_v2's
% tables/t2_graph.tex). Numbers from numbers.tex (data/AMAZON/meta.json, audit.json, label files;
% the split rebuilt and checked by make_numbers.py).
\begin{widetable}[htbp]
  \centering
  \small
  \caption{Đồ thị và phép chia. Mọi đồ thị mà các mô hình thấy chỉ chứa các cạnh huấn luyện.}
  \label{tab:graph}
  \begin{tabularx}{\linewidth}{@{}lrL@{}}
    \toprule
    Kiểu nút & Số nút & Đặc trưng \\
    \midrule
    Sản phẩm (P) & \val{after-products} & vectơ LSA \val{title-dim} chiều của tiêu đề (khác 0 với
      \val{title-coverage-pct}\%) và nhóm sản phẩm: sách \val{group-book-pct}\%, nhạc
      \val{group-music-pct}\%, DVD \val{group-dvd-pct}\%, video \val{group-video-pct}\% \\
    Danh mục (C) & \val{after-categories} & vectơ multi-hot \val{cat-vocab} chiều của các token trong
      tên (toàn số 0 với \val{after-zerocat-pct}\%) \\
    Người dùng (U) & \val{after-users} & vectơ nhị phân \val{user-dim} chiều sinh từ ID bằng hàm băm. Trong tìm kiếm, HAN và
      metapath2vec thay bằng một embedding tự do \val{search-hidden} chiều \\
    \midrule
    Quan hệ & Số cạnh & \\
    \midrule
    P--P, mua kèm & \val{after-alsobought} & cặp sản phẩm (\val{rel-pp-rows} dòng có hướng), đã
      đối xứng hóa \\
    C--C, cha và con & \val{rel-cc} & \\
    P--C, danh mục lá & \val{rel-pc} & trung vị \val{after-leafcats-med} danh mục mỗi sản phẩm \\
    U--P, nhận xét tích cực & \val{rel-up} & chia ngẫu nhiên một lần: \val{split-train} cạnh huấn
      luyện, \val{split-val} cạnh xác thực, \val{split-test} cạnh thử nghiệm (seed của phép
      chia~\val{split-seed}) \\
    \midrule
    \multicolumn{3}{@{}>{\raggedright\arraybackslash}p{\linewidth}@{}}{Số cạnh huấn luyện mỗi
      người dùng: trung vị \val{train-deg-med}, trung bình \val{train-deg-mean}, lớn nhất
      \val{train-deg-max}. Cạnh thử nghiệm theo số cạnh huấn luyện của người dùng:
      \val{testdeg-d1} (1--5), \val{testdeg-d2} (6--20), \val{testdeg-d3} (từ 21 trở lên) và
      \val{testdeg-d0} (không có). Mỗi cạnh xác thực và thử nghiệm được xếp hạng so với
      \val{split-negatives} mẫu âm lấy mẫu và so với mọi sản phẩm.} \\
    \bottomrule
  \end{tabularx}
\end{widetable}


\section{Thay đổi so với các bài báo gốc}
\label{app:changes}

Bảng~\ref{tab:changes} liệt kê những gì nghiên cứu này thay đổi trong ba phương pháp gốc, kèm lý do.
Phần lớn các thay đổi xuất phát từ bài toán. LMSPS và HAN được xây dựng cho bài toán phân loại nút.
Dự đoán liên kết cần một điểm cho cặp người dùng--sản phẩm, một hàm mất mát trên các cặp như vậy
và, trong tìm kiếm, sự cẩn trọng để một cạnh huấn luyện không chạm tới chính đặc trưng của nó. Các
thay đổi khác làm cho phép so sánh công bằng: mọi tập meta-path đều gặp cùng một phía người dùng,
cùng các thiết lập và cùng quy tắc dừng, trong HAN cũng như trong metapath2vec. Phần còn lại làm cho
quyết định của tìm kiếm tái lập được: tìm kiếm tính điểm mọi tập của nhóm ứng viên cuối mà không có
nhiễu, rồi nhiều lần chạy được kết hợp thành một lựa chọn. Thiết lập của các bài báo gốc được ghi
như chúng báo cáo, còn thiết lập của chúng tôi là thiết lập của các lần chạy (Mục~\ref{app:settings}).

% t6_changes.tex -- T6 (appendix): what this study changes in the three base methods (paper_v2_vn,
% translated in L3 from paper_v2's tables/t6_changes.tex). The base papers' own settings are typed
% from the papers, here in the Vietnamese number format (S5 of glossary_vn.tex); this study's
% settings are \val{} from numbers.tex (run records).
\begin{widetable}[htbp]
  \centering
  \footnotesize
  \renewcommand{\arraystretch}{1.15}
  \caption{Những gì nghiên cứu này thay đổi so với LMSPS \citep{li2024longrange}, HAN
    \citep{wang2019heterogeneous} và metapath2vec \citep{dong2017metapath2vec}, kèm lý do.}
  \label{tab:changes}
  \begin{tabularx}{\linewidth}{@{}>{\raggedright\arraybackslash}p{0.27\linewidth}
      >{\raggedright\arraybackslash}p{0.34\linewidth}L@{}}
    \toprule
    Trong bài báo gốc & Trong nghiên cứu này & Lý do \\
    \midrule
    \multicolumn{3}{@{}l}{\emph{LMSPS}} \\
    phân loại nút, entropy chéo & dự đoán liên kết từ người dùng đến sản phẩm: BPR trên tích vô
      hướng của embedding sản phẩm và một embedding người dùng tự do \val{search-hidden} chiều & bài toán
      của nghiên cứu này \\
    mọi meta-path từ kiểu đích đến một kiểu bất kỳ, kể cả đặc trưng gốc & \val{search-cands}
      meta-path từ sản phẩm đến sản phẩm dài 1--\val{search-hops} bước, mỗi meta-path gộp với
      meta-path ngược của nó, không có cạnh tự nối & mỗi meta-path phải cho HAN một đồ thị sản
      phẩm \\
    đặc trưng meta-path tính từ toàn bộ đồ thị & $\alpha$ học trên các cạnh xác thực. Bước cập
      nhật trọng số đọc đặc trưng meta-path qua người dùng được dựng lại khi bỏ lần lượt từng phần
      tư số cạnh huấn luyện & mỗi cạnh huấn luyện cũng là một cạnh của đồ thị \\
    $M = 30$ meta-path lấy mẫu từ nhóm ứng viên $V = 60$, 200 epoch, trong đó 20 epoch khởi động
      (mã công bố không có) & $M = 4$ hoặc 6 từ $V = \val{search-pool}$, \val{search-epochs}
      epoch, trong đó \val{search-warmup} epoch khởi động, dropout \val{search-dropout} & các kích
      thước tập của nghiên cứu, để mọi tập của nhóm ứng viên cuối đều tính điểm được \\
    200 tập lấy mẫu từ nhóm ứng viên cuối, mỗi tập tính điểm với nhiễu Gumbel mới & cả
      \val{search-sets} tập được tính điểm một lần, không có nhiễu & nhiễu làm điểm của một tập
      thay đổi nhiều hơn chênh lệch giữa các tập \\
    một lần tìm kiếm & \val{search-runs} lần tìm kiếm cho mỗi kích thước tập, kết hợp bằng điểm
      trung bình với quy tắc một sai số chuẩn và bằng bỏ phiếu được kiểm tra so với đối chứng xáo
      trộn & một lần chạy không xếp hạng được các tập tốt nhất của nó \\
    mạng riêng của LMSPS được huấn luyện lại trên các meta-path đã chọn & không huấn luyện lại:
      các tập được chuyển cho HAN và metapath2vec & câu hỏi là các meta-path có giúp HAN hay
      không \\
    \midrule
    \multicolumn{3}{@{}l}{\emph{HAN}} \\
    phân loại nút, entropy chéo & dự đoán liên kết, BPR. HAN tạo embedding cho sản phẩm, mỗi
      người dùng là một embedding tự do \val{han-set-hidden} chiều & mọi tập meta-path có cùng phía
      người dùng \\
    mọi láng giềng theo meta-path & \val{han-set-topk} láng giềng mạnh nhất sau mỗi bước,
      đối xứng, chỉ cạnh huấn luyện, một cạnh tự nối & meta-path dài không cắt tỉa đến được hầu hết
      sản phẩm \\
    tầng chú ý của mỗi meta-path có khởi tạo ngẫu nhiên riêng. Trọng số meta-path được học ngay từ
      epoch đầu & cùng khởi tạo cho mọi meta-path. Trọng số giữ bằng nhau trong
      \val{han-set-beta-warmup} epoch & nếu không, một meta-path chiếm toàn bộ trọng số chỉ sau vài
      epoch \\
    tốc độ học 0,005, L2 0,001, dropout 0,6, 8 đầu chú ý 8 chiều, patience 100 & tốc độ học
      \val{han-set-lr}, không L2, dropout \val{han-set-feat-drop}, \val{han-set-num-heads} đầu chú
      ý \val{han-set-head-dim} chiều, patience \val{han-set-patience}, tối đa \val{han-set-stage}
      epoch & vừa với một GPU \val{env-gpu-mem}~GB \\
    đặc trưng túi từ, 2--3 meta-path ngắn đối xứng & vectơ LSA của tiêu đề và nhóm sản phẩm,
      \val{sets-size-min}--\val{sets-size-max} meta-path dài tối đa \val{sets-hops-max} bước & dữ
      liệu và các tập được so sánh \\
    \midrule
    \multicolumn{3}{@{}l}{\emph{metapath2vec}} \\
    mỗi lần chạy một lược đồ meta-path đối xứng & mọi meta-path của một tập được đi cùng nhau &
      mỗi tập một embedding, như với HAN \\
    mẫu âm theo phân phối tần suất nút & \val{mtv-set-negatives} mẫu âm, lấy mẫu đều trong cùng kiểu nút
      & đơn giản hơn \\
    1.000 bước đi mỗi nút, độ dài 100, cửa sổ 7, 128 chiều & \val{mtv-set-walks-per-node} bước đi
      cho mỗi nút và mỗi meta-path, độ dài \val{mtv-set-walk-length}, cửa sổ \val{mtv-set-window},
      \val{mtv-set-dim} chiều & chi phí, bằng kích thước embedding của HAN \\
    \midrule
    \multicolumn{3}{@{}l}{\emph{GAT, làm mô hình so sánh}} \\
    trong bài báo HAN, chạy trên từng meta-path riêng lẻ & một GAT không phân biệt kiểu trên
      mọi quan hệ cùng lúc, không dùng meta-path & mô hình so sánh đồng nhất của nghiên cứu
      này \\
    \bottomrule
  \end{tabularx}
\end{widetable}


\section{Chi tiết về tìm kiếm}
\label{app:search}

\paragraph{Mức đồng thuận giữa các lần chạy.} Các lần chạy tính điểm những tập khác nhau, vì nhóm
ứng viên cuối của chúng khác nhau, nhưng chúng xếp hạng giống nhau những tập mà chúng cùng có: điểm
của hai lần chạy trên các tập chung có hệ số tương quan Spearman trung vị \val{search-agree-4} với
bốn meta-path và \val{search-agree-6} với sáu, trong khi điểm của một tập dao động
\val{search-noise-4} và \val{search-noise-6} (độ lệch chuẩn) từ lần chạy này sang lần chạy khác
(Mục~\ref*{P-sec:fleet}).

\paragraph{Bỏ phiếu so với đối chứng.} Bỏ phiếu chọn ra \val{search-vote-4} tập khác nhau với bốn
meta-path, so với trung vị \val{search-null-med-4} (khoảng 90\%
\val{search-null-lo-4}--\val{search-null-hi-4}) trong \val{search-null-fleets} loạt chạy có điểm bị
xáo trộn trong từng lần chạy. Với sáu meta-path, bỏ phiếu chọn ra \val{search-vote-6} tập,
so với \val{search-null-med-6} (\val{search-null-lo-6}--\val{search-null-hi-6}). Hình~\ref*{P-fig:ties}b thể
hiện cả hai.

\paragraph{Thứ hạng của các meta-path.} Qua \val{search-runs} lần chạy của mỗi kích thước, hạng
$\alpha$ trung bình của \code{PUPUP} và \code{PUPPUP} là \val{rank-pupup-4} và \val{rank-puppup-4}
trong \val{search-cands} với bốn meta-path, còn với sáu là \val{rank-pupup-6} và \val{rank-puppup-6}.
Các meta-path danh mục dày đặc mà các tập thủ công dùng xếp hạng thấp: \code{PCP} ở
\val{rank-pcp-4} và \val{rank-pcp-6}, \code{PCCCP} ở \val{rank-pcccp-4} và \val{rank-pcccp-6}. Lấy
trung bình trên cả hai kích thước, các meta-path danh mục gần rỗng xếp hạng \val{kind-empty-rank}
và các meta-path dày đặc \val{kind-cat-rank} (Hình~\ref*{P-fig:weights}).

\paragraph{Dung sai của tìm kiếm.} Với $\delta = \val{search-delta-alt}$ thay vì
\val{search-delta}, bỏ phiếu chọn ra \val{search-vote-alt-4} và \val{search-vote-alt-6} tập khác
nhau, so với trung vị trên điểm xáo trộn là \val{search-null-med-alt-4} và
\val{search-null-med-alt-6}, nên các lần chạy vẫn đồng thuận hơn mức ngẫu nhiên. Các điểm trung
bình không phụ thuộc vào $\delta$, nên \code{a1} và \code{a2} cũng vậy.

\section{Các tập chọn ngẫu nhiên và~các~thiết~lập}
\label{app:settings}

Với một sản phẩm liên quan cho mỗi lần xếp hạng, Recall@10 bằng Hit@10 và Precision@10 bằng
Hit@10$/10$, nên Mục~\ref*{P-sec:setup} không báo cáo hai độ đo này. Bảng~\ref{tab:random} liệt kê mười
tập chọn ngẫu nhiên của Mục~\ref*{P-sec:setup}. Bảng~\ref{tab:settings} nêu mọi thiết lập của tìm
kiếm và của các mô hình. Mỗi thiết lập được cố định trước các lần chạy và dùng cho mọi tập meta-path (Mục~\ref*{P-sec:setup}). Các giá trị được đọc từ bản ghi và nhật ký của các lần chạy: của tìm
kiếm từ \val{search-runs-total} bản ghi (vốn giống nhau ở mọi thiết lập trừ seed và kích thước
tập), còn của các mô hình từ các lệnh và phần mô tả mô hình trong nhật ký của chúng. Mọi lần chạy được báo
cáo đều dừng sớm: HAN ở epoch \val{han-stop-min}--\val{han-stop-max} trong tối đa
\val{han-set-stage} epoch cho phép, GAT ở \val{gat-stop-min}--\val{gat-stop-max} (một seed sau khi
chạy lại với mức trần cao hơn) và metapath2vec sau \val{mtv-stop-min}--\val{mtv-stop-max} lượt.

% t3b_random.tex -- the ten random draws, in the supplementary document (paper_v2_vn, translated
% in L3 from paper_v2's tables/t3b_random.tex). The paper's table of the sets (t3_sets.tex) lists
% them too since N31. Paths from numbers.tex (HAN repository, arch.py).
\begin{widetable}[htbp]
  \centering
  \small
  \caption{Mười tập chọn ngẫu nhiên, chọn đều và không lặp lại từ \val{search-cands} ứng viên của
    tìm kiếm (Mục~\ref*{P-sec:setup}).}
  \label{tab:random}
  \begin{tabular}{@{}ll@{}}
    \toprule
    Lần chọn & Meta-path \\
    \midrule
    4 meta-path, lần chọn 1 & \texttt{\val{set-r4d1}} \\
    4 meta-path, lần chọn 2 & \texttt{\val{set-r4d2}} \\
    4 meta-path, lần chọn 3 & \texttt{\val{set-r4d3}} \\
    4 meta-path, lần chọn 4 & \texttt{\val{set-r4d4}} \\
    4 meta-path, lần chọn 5 & \texttt{\val{set-r4d5}} \\
    \addlinespace
    6 meta-path, lần chọn 1 & \texttt{\val{set-r6d1}} \\
    6 meta-path, lần chọn 2 & \texttt{\val{set-r6d2}} \\
    6 meta-path, lần chọn 3 & \texttt{\val{set-r6d3}} \\
    6 meta-path, lần chọn 4 & \texttt{\val{set-r6d4}} \\
    6 meta-path, lần chọn 5 & \texttt{\val{set-r6d5}} \\
    \bottomrule
  \end{tabular}
\end{widetable}

% t8_settings.tex -- the settings of the search and of every model (appendix; paper_v2_vn,
% translated in L3 from paper_v2's tables/t8_settings.tex). Values from numbers.tex: the search's
% from its 60 run records, which make_numbers.py checks are alike apart from the seed and the set
% size; the models' from the commands and model prints in their logs.
\begin{widetable}[htbp]
  \centering
  \footnotesize
  \renewcommand{\arraystretch}{1.0}
  \caption{Mọi thiết lập của tìm kiếm và của các mô hình, cố định trước các lần chạy và dùng cho
    mọi tập meta-path.}
  \label{tab:settings}
  \begin{tabularx}{\linewidth}{@{}>{\raggedright\arraybackslash}p{0.25\linewidth}L@{}}
    \toprule
    Thiết lập & Giá trị \\
    \midrule
    \multicolumn{2}{@{}l}{\emph{Tìm kiếm (LMSPS cho dự đoán liên kết)}} \\
    Ứng viên & \val{search-cands-raw} meta-path từ sản phẩm đến sản phẩm dài 1 đến \val{search-hops}
      bước, còn \val{search-cands} khi mỗi meta-path được gộp với meta-path ngược của nó. Không có
      cạnh tự nối \\
    Tập và nhóm ứng viên & $M = 4$ hoặc 6 meta-path. Nhóm ứng viên cuối $V = \val{search-pool}$,
      nên cả \val{search-sets} tập gồm $M$ meta-path đều được tính điểm \\
    Lịch huấn luyện & \val{search-epochs} epoch, \val{search-warmup} epoch đầu với toàn bộ nhóm
      ứng viên. Nhiệt độ Gumbel giảm từ \val{search-tau-max} xuống \val{search-tau-min} \\
    Mạng & mỗi meta-path một ánh xạ tuyến tính, PReLU, dropout \val{search-dropout} (dropout đầu
      vào \val{search-input-drop}). Embedding sản phẩm và người dùng \val{search-hidden} chiều \\
    Bộ tối ưu & Adam: tốc độ học \val{search-lr} cho trọng số và \val{search-alr} cho $\alpha$,
      suy giảm trọng số \val{search-weight-decay}. Lô gồm \val{search-batch} cạnh \\
    Chống rò rỉ & $\alpha$ học trên các cạnh xác thực. Bước cập nhật trọng số đọc đặc trưng meta-path qua người dùng từ \val{search-folds} phiên bản, mỗi phiên bản dựng khi bỏ một phần tư số
      cạnh huấn luyện \\
    Tính điểm & mỗi tập một lần, không có nhiễu Gumbel, theo BPR xác thực so với
      \val{split-negatives} mẫu âm cố định. Nhóm đồng hạng trong phạm vi
      $\delta = \val{search-delta}$ \\
    Lần chạy & \val{search-runs} cho mỗi kích thước tập, seed 1 đến \val{search-runs}. Kernel GPU
      tất định. Đặc trưng lan truyền lưu dạng số thực 16 bit trên CPU \\
    \midrule
    \multicolumn{2}{@{}l}{\emph{HAN}} \\
    Đồ thị & theo từng meta-path, \val{han-set-topk} láng giềng mạnh nhất của mỗi sản phẩm theo số
      đường đi sau mỗi bước, bằng điểm thì xếp ngẫu nhiên. Nhị phân, đối xứng, mỗi sản phẩm một
      cạnh tự nối. Chỉ cạnh huấn luyện \\
    Bộ mã hóa & đặc trưng sản phẩm \val{han-set-in-dim} chiều. Mỗi meta-path một tầng chú ý đồ
      thị gồm \val{han-set-num-heads} đầu chú ý \val{han-set-head-dim} chiều, mọi meta-path bắt đầu
      từ cùng trọng số khởi tạo \\
    Trọng số meta-path & chú ý mức ngữ nghĩa với vectơ chú ý \val{han-set-sem-dim} chiều và nhiệt
      độ softmax \val{han-set-semantic-tau}. Trọng số giữ bằng nhau trong
      \val{han-set-beta-warmup} epoch đầu \\
    Người dùng & mỗi người dùng một embedding tự do \val{han-set-hidden} chiều \\
    Huấn luyện & dropout \val{han-set-feat-drop} trên đặc trưng và \val{han-set-attn-drop} trên
      chú ý. Adam, tốc độ học \val{han-set-lr}, suy giảm trọng số \val{han-set-weight-decay}. Lô
      gồm \val{han-set-batch-size} cạnh huấn luyện, mỗi cạnh một mẫu âm mới \\
    Dừng & tối đa \val{han-set-stage} epoch. Dừng khi BPR xác thực không cải thiện quá
      \val{han-set-min-delta} trong \val{han-set-patience} epoch. Thử nghiệm điểm lưu tốt nhất \\
    \midrule
    \multicolumn{2}{@{}l}{\emph{metapath2vec}} \\
    Bước đi & với mỗi sản phẩm và meta-path, \val{mtv-set-walks-per-node} bước đi dài
      \val{mtv-set-walk-length}. Mọi meta-path của một tập được đi cùng nhau, mỗi meta-path như một
      khuôn mẫu lặp lại. Chỉ cạnh huấn luyện \\
    Skip-gram & \val{mtv-set-dim} chiều, cửa sổ \val{mtv-set-window}, \val{mtv-set-negatives} mẫu
      âm mỗi cặp, lấy mẫu đều từ kiểu của nút ngữ cảnh \\
    Huấn luyện & Adam thưa (sparse Adam), tốc độ học \val{mtv-set-lr}. Lô gồm
      \val{mtv-set-batch-walks} bước đi \\
    Dừng & tối đa \val{mtv-set-epochs} lượt qua các bước đi. Kiểm tra BPR xác thực
      \val{mtv-set-evals-per-epoch} lần mỗi lượt. Dừng sau \val{mtv-set-patience} lần kiểm tra
      không cải thiện quá \val{mtv-set-min-delta} \\
    \midrule
    \multicolumn{2}{@{}l}{\emph{GAT}} \\
    Bộ mã hóa & mỗi kiểu nút một phép chiếu tuyến tính về \val{gat-set-in-dim} chiều, sau đó
      \val{gat-set-layers} tầng gồm \val{gat-set-num-heads} đầu chú ý (tổng cộng
      \val{gat-set-hidden} chiều) trên hợp của mọi quan hệ, ngoài ra bỏ qua kiểu nút \\
    Huấn luyện & dropout \val{gat-set-feat-drop} trên đặc trưng và \val{gat-set-attn-drop} trên
      chú ý. Adam, tốc độ học \val{gat-set-lr}, suy giảm trọng số \val{gat-set-weight-decay}. Lô
      gồm \val{gat-set-batch-size} cạnh huấn luyện, các cạnh dương của lô được bỏ khỏi đồ thị mà
      bước đó chạy trên \\
    Dừng & như HAN, tối đa \val{gat-cap} epoch. Seed~\val{gat-rerun-seed} chạm mức trần đó và được
      chạy lại với \val{gat-rerun-cap}, dừng ở epoch \val{gat-rerun-stop} \\
    \midrule
    \multicolumn{2}{@{}l}{\emph{Mô hình phổ biến}} \\
    Điểm & số người dùng huấn luyện của sản phẩm, bằng điểm thì xếp ngẫu nhiên. Một bảng xếp hạng
      chung cho mọi người dùng \\
    \midrule
    \multicolumn{2}{@{}l}{\emph{Đánh giá}} \\
    Phép chia & 80/10/10 ngẫu nhiên, seed của phép chia~\val{split-seed}. Mỗi cạnh xác thực
      và thử nghiệm có \val{split-negatives} mẫu âm, lấy mẫu một lần \\
    Độ đo & Hit@10 lấy mẫu (cố định từ trước) và NDCG@10. Khi xếp hạng toàn bộ, Hit@10 và NDCG@10
      theo cạnh và Recall@20 theo người dùng \\
    Seed & 1 đến 3 cho so sánh chính và GAT, 1 cho các tập khác \\
    \bottomrule
  \end{tabularx}
\end{widetable}


\section{Kết quả đầy đủ}
\label{app:full}

Bảng~\ref{tab:full-han} và~\ref{tab:full-mtv} nêu mọi độ đo cho mọi tập meta-path ở seed~1, seed mà
mọi tập đều được chạy. Bảng~\ref*{P-tab:main} nêu trung bình ba seed của so sánh chính. Qua các seed,
Hit@10 thử nghiệm của các lần chạy chính dao động nhiều nhất \val{main-hit-sd-max} (độ lệch chuẩn,
Mục~\ref*{P-sec:setup}), nên các chênh lệch nhỏ trong các bảng này nằm trong nhiễu giữa các seed.

Bảng~\ref{tab:groups} nêu Hit@10 thử nghiệm của mọi mô hình theo số cạnh huấn luyện của người dùng
(Mục~\ref*{P-sec:res-models}). Phần còn lại của mục này nêu những kết quả mà
Mục~\ref*{P-sec:results} không trình bày.

\paragraph{Nhóm người dùng.} \code{a1} của HAN kém tập thủ công cùng kích thước ở mọi nhóm người
dùng, với chênh lệch \val{han-4-hit-d1-diff}, \val{han-4-hit-d2-diff} và \val{han-4-hit-d3-diff} ở
bốn meta-path và \val{han-6-hit-d1-diff}, \val{han-6-hit-d2-diff} và \val{han-6-hit-d3-diff} ở sáu,
cho người dùng có 1--5, 6--20 và từ 21 cạnh huấn luyện trở lên. Mức dẫn trước của GAT so với tập
chọn theo xác thực của HAN là \val{gat-lead-d1}, \val{gat-lead-d2} và \val{gat-lead-d3} trong cùng
các nhóm.

\paragraph{Thêm meta-path.} Thêm meta-path giúp HAN ở mọi loại tập: các tập tìm được đạt
\val{han-searched6-mean-s1} với sáu meta-path so với \val{han-searched4-mean-s1} với bốn, còn các
tập chọn ngẫu nhiên đạt \val{han-r6-hit-mean} so với \val{han-r4-hit-mean}.

\paragraph{Trọng số và mức ưa thích meta-path.} HAN gán cho \code{PUPPUP}, \code{PUPUP} và
\code{PUP} trọng số bằng \val{weight-puppup}, \val{weight-pupup} và \val{weight-pup} lần mức chia
đều, đồng thời loại \code{PPCCPP} ở cả \val{uses-ppccpp} tập chứa nó. Trên \val{pref-cand-n} tập lấy từ các
ứng viên, hạng $\alpha$ trung bình của các meta-path trong một tập cũng có hệ số tương quan Spearman
\val{pref-cand-val-rho} với BPR xác thực của HAN (\val{pref-cand-val-pt}). Riêng trong mười tập chọn
ngẫu nhiên, tương quan của nó với Hit@10 thử nghiệm là yếu (\val{pref-random-rho},
\val{pref-random-pt}), còn trên cả 24 tập là \val{pref-all-rho} (\val{pref-all-pt}): quan hệ này dựa
trên các tập tìm được.

\paragraph{HAN và metapath2vec.} Hai mô hình ưa thích các tập khác nhau: trên 23 tập có meta-path
qua người dùng, giá trị Hit@10 của chúng có hệ số tương quan Spearman chỉ \val{han-mtv-rho}. Chúng
cũng gợi ý những sản phẩm rất khác nhau. Trong top 10 xếp hạng toàn bộ của \val{bg-users} người dùng
thử nghiệm, các sản phẩm của HAN có trung vị \val{bg-han-deg-med} người dùng huấn luyện
(\val{bg-han-deg-min}--\val{bg-han-deg-max} giữa các tập), của metapath2vec là \val{bg-mtv-deg-med}
và của GAT là \val{bg-gat-deg}, so với \val{bg-all-deg} của một sản phẩm trung bình. Trong các gợi ý
của HAN, \val{bg-han-top1-med}\% đến từ 1\% sản phẩm được mua nhiều nhất, so với
\val{bg-gat-top1}\% của GAT và nhiều nhất \val{bg-mtv-top1-max}\% của metapath2vec. Theo nhóm sản phẩm, HAN
dẫn trước metapath2vec \val{bg-han-mtv-dvd} với DVD và \val{bg-han-mtv-video} với video nhưng kém nó
với nhạc (\val{bg-han-mtv-music}, trung vị trên 23 tập). Hai mô hình đổi chỗ cho nhau giữa xếp hạng
lấy mẫu và xếp hạng toàn bộ (Mục~\ref*{P-sec:res-models}): các sản phẩm phổ biến của HAN dễ dàng thắng
\val{split-negatives} mẫu âm ngẫu nhiên, trong khi các sản phẩm đuôi dài của metapath2vec làm tốt
hơn một cách tương đối khi so với mọi sản phẩm, đúng loại bất đồng mà \citet{krichene2022sampled}
đã cảnh báo.

% t5_groups.tex -- T5: every model, overall and by the user's number of training edges
% (paper_v2_vn, translated in L3 from paper_v2's tables/t5_groups.tex; results.csv and the run
% logs, via numbers.tex).
\begin{widetable}[htbp]
  \centering
  \small
  \caption{Mọi mô hình, trên toàn bộ và theo số cạnh huấn luyện của người dùng: Hit@10 trên tập
    thử nghiệm, trung bình qua các seed đã chạy. GAT, các tập thủ công và \code{a1} chạy với 3 seed.
    Tập chọn theo xác thực của HAN (một tập ngẫu nhiên) và mô hình phổ biến chạy với 1 seed.
    Kích thước nhóm là số cạnh thử nghiệm.}
  \label{tab:groups}
  \begin{tabular*}{\linewidth}{@{\extracolsep{\fill}}lrrrrr@{}}
    \toprule
    & \multicolumn{2}{c}{Mọi cạnh thử nghiệm} & \multicolumn{3}{c}{Số cạnh huấn luyện của người
      dùng} \\
    \cmidrule(lr){2-3}\cmidrule(l){4-6}
    Mô hình và tập & Lấy mẫu & Toàn bộ & 1--5 & 6--20 & từ 21 trở lên \\
    & & & (\val{testdeg-d1}) & (\val{testdeg-d2}) & (\val{testdeg-d3}) \\
    \midrule
    GAT, không dùng meta-path & \val{gat-hit-mean} & \val{gat-fhit-mean} & \val{gat-hit-d1}
      & \val{gat-hit-d2} & \val{gat-hit-d3} \\
    HAN, tập chọn theo xác thực (6 meta-path, lần chọn 2) & \val{han-r6d2-hit-s1}
      & \val{han-r6d2-fhit-s1} & \val{han-r6d2-hit-d1} & \val{han-r6d2-hit-d2}
      & \val{han-r6d2-hit-d3} \\
    HAN, M6 & \val{han-m6-hit-mean} & \val{han-m6-fhit-mean} & \val{han-m6-hit-d1}
      & \val{han-m6-hit-d2} & \val{han-m6-hit-d3} \\
    HAN, \code{a1}, 6 meta-path & \val{han-f6a1-hit-mean} & \val{han-f6a1-fhit-mean}
      & \val{han-f6a1-hit-d1} & \val{han-f6a1-hit-d2} & \val{han-f6a1-hit-d3} \\
    metapath2vec, M5 (tập chọn theo xác thực) & \val{mtv-m5-hit-mean} & \val{mtv-m5-fhit-mean}
      & \val{mtv-m5-hit-d1} & \val{mtv-m5-hit-d2} & \val{mtv-m5-hit-d3} \\
    metapath2vec, \code{a1}, 4 meta-path & \val{mtv-f4a1-hit-mean} & \val{mtv-f4a1-fhit-mean}
      & \val{mtv-f4a1-hit-d1} & \val{mtv-f4a1-hit-d2} & \val{mtv-f4a1-hit-d3} \\
    Phổ biến & \val{pop-hit} & \val{pop-fhit} & \val{pop-hit-d1} & \val{pop-hit-d2}
      & \val{pop-hit-d3} \\
    \bottomrule
  \end{tabular*}
\end{widetable}

% t9_full.tex -- every path set in HAN and in metapath2vec, every metric, seed 1 (supplementary
% document; paper_v2_vn, translated in L3 from paper_v2's tables/t9_full.tex). Values from
% numbers.tex (results.csv): <model>-<set>-<metric>-s1; GAT and popularity as in T4/T5.
% \fullrow{label}{model tag}{set}: one row of six metrics.
\newcommand{\fullrow}[3]{#1 & \val{#2-#3-hit-s1} & \val{#2-#3-ndcg-s1} & \val{#2-#3-fhit-s1}
  & \val{#2-#3-fndcg-s1} & \val{#2-#3-frec-s1} & \val{#2-#3-val-s1} \\}
% \fullbody{model tag}: the 24 sets, by family, as in Table 3.
\newcommand{\fullbody}[1]{%
  \multicolumn{7}{@{}l}{\emph{Tìm được}} \\
  \fullrow{\code{a1}, 4 meta-path}{#1}{f4a1}
  \fullrow{\code{a2}, 4 meta-path}{#1}{f4a2}
  \fullrow{\code{c}, 4 meta-path}{#1}{f4c}
  \fullrow{\code{a1}, 6 meta-path}{#1}{f6a1}
  \fullrow{\code{a2}, 6 meta-path}{#1}{f6a2}
  \fullrow{\code{c}, 6 meta-path}{#1}{f6c}
  \midrule
  \multicolumn{7}{@{}l}{\emph{Thủ công}} \\
  \fullrow{M1}{#1}{m1}
  \fullrow{M2}{#1}{m2}
  \fullrow{M3}{#1}{m3}
  \fullrow{M4}{#1}{m4}
  \fullrow{M5}{#1}{m5}
  \fullrow{M6}{#1}{m6}
  \fullrow{M7}{#1}{m7}
  \fullrow{M8}{#1}{m8}
  \midrule
  \multicolumn{7}{@{}l}{\emph{Ngẫu nhiên}} \\
  \fullrow{4 meta-path, lần chọn 1}{#1}{r4d1}
  \fullrow{4 meta-path, lần chọn 2}{#1}{r4d2}
  \fullrow{4 meta-path, lần chọn 3}{#1}{r4d3}
  \fullrow{4 meta-path, lần chọn 4}{#1}{r4d4}
  \fullrow{4 meta-path, lần chọn 5}{#1}{r4d5}
  \fullrow{6 meta-path, lần chọn 1}{#1}{r6d1}
  \fullrow{6 meta-path, lần chọn 2}{#1}{r6d2}
  \fullrow{6 meta-path, lần chọn 3}{#1}{r6d3}
  \fullrow{6 meta-path, lần chọn 4}{#1}{r6d4}
  \fullrow{6 meta-path, lần chọn 5}{#1}{r6d5}}
\newcommand{\fullhead}{%
  \toprule
  & \multicolumn{2}{c}{Lấy mẫu} & \multicolumn{3}{c}{Xếp hạng toàn bộ} & \\
  \cmidrule(lr){2-3}\cmidrule(lr){4-6}
  Tập & Hit@10 & NDCG@10 & Hit@10 & NDCG@10 & Recall@20 & BPR xác thực \\
  \midrule}

\begin{widetable}[htbp]
  \centering
  \small
  \caption{HAN trên mọi tập meta-path, seed~1 (Bảng~\ref*{P-tab:main} cho trung bình ba seed của so
    sánh chính). Lấy mẫu: mỗi cạnh thử nghiệm được xếp hạng so với \val{split-negatives} mẫu âm.
    Xếp hạng toàn bộ: so với mọi sản phẩm mà người dùng không có cạnh nối tới, với Recall@20 lấy
    trung bình theo người dùng. BPR xác thực: giá trị hàm mất mát mà việc dừng sớm và tập chọn
    theo xác thực sử dụng, càng thấp càng tốt. GAT: trung bình ba seed. Mô hình phổ biến không có
    giá trị mất mát xác thực.}
  \label{tab:full-han}
  \begin{tabular*}{\linewidth}{@{\extracolsep{\fill}}lllllll@{}}
    \fullhead
    \fullbody{han}
    \midrule
    \multicolumn{7}{@{}l}{\emph{Không dùng meta-path}} \\
    GAT & \val{gat-hit-mean} & \val{gat-ndcg-mean} & \val{gat-fhit-mean} & \val{gat-fndcg-mean}
      & \val{gat-frec-mean} & \val{gat-val-mean} \\
    Phổ biến & \val{pop-hit} & \val{pop-ndcg} & \val{pop-fhit} & \val{pop-fndcg} & \val{pop-frec}
      & -- \\
    \bottomrule
  \end{tabular*}
\end{widetable}

\begin{widetable}[htbp]
  \centering
  \small
  \caption{metapath2vec trên mọi tập meta-path, seed~1. Các cột như ở Bảng~\ref{tab:full-han}. M4
    không có meta-path nào qua người dùng, nên các embedding người dùng của nó không được huấn luyện.}
  \label{tab:full-mtv}
  \begin{tabular*}{\linewidth}{@{\extracolsep{\fill}}lllllll@{}}
    \fullhead
    \fullbody{mtv}
    \bottomrule
  \end{tabular*}
\end{widetable}


\section{Khả năng tái lập}
\label{app:repro}

\paragraph{Mã nguồn.} Quy trình xử lý dữ liệu và tìm kiếm nằm trong kho mã LMSPS, các mô hình và
phần đánh giá trong kho mã HAN (mục Mã nguồn và dữ liệu của bài báo). Các kết quả trong bài báo đến
từ commit
\texttt{\val{repo-lmsps-commit}} của kho thứ nhất (nhánh \texttt{\val{repo-lmsps-branch}}) và
commit \texttt{\val{repo-han-commit}} của kho thứ hai (nhánh \texttt{\val{repo-han-branch}}).

\paragraph{Dữ liệu.} Amazon-Meta lấy từ SNAP (Mục~\ref*{P-sec:data}). Script \code{generate_data.bat}
trong kho mã LMSPS (thư mục \code{hgb}) trích xuất và làm sạch tệp thô, ghi các tệp đồ thị và rà
soát kết quả, với các thiết lập ở Mục~\ref*{P-sec:data}. Mỗi lần chạy ghi lại một dấu vân tay
(fingerprint) của các cạnh huấn luyện trong phép chia của nó, \texttt{\val{split-fingerprint}} với
seed của phép chia~\val{split-seed}. Script \code{check_handover.py} của kho mã HAN xác nhận rằng dữ
liệu của nó chính là dữ liệu mà tìm kiếm đã chạy trên.

\paragraph{Môi trường.} Một GPU NVIDIA \val{env-gpu} (\val{env-gpu-mem}\,GB) chạy Windows, với
Python \val{env-python}, PyTorch \val{env-torch} xây dựng cho CUDA \val{env-cuda} và NumPy
\val{env-numpy}. Tìm kiếm còn dùng DGL \val{env-dgl} cho bước lan truyền đặc trưng một lần trên
CPU. Tệp \code{requirements.txt} của mỗi kho mã nêu phiên bản các gói của nó. Các script của cả
hai kho chạy trong một môi trường conda. Các lần chạy mô hình mất khoảng \val{models-hours}~giờ~GPU:
\val{han-hours} cho \val{han-runs} lần chạy HAN, \val{mtv-hours} cho metapath2vec và
\val{gat-hours} cho GAT.

\paragraph{Các lệnh.} Trong kho mã LMSPS (thư mục \code{hgb}), \code{run_fleet_v3.bat pilot} trước
hết kiểm tra rằng hai lần chạy với cùng seed cho kết quả giống nhau. Sau đó,
\code{run_fleet_v3.bat final4} và \code{final6} chạy \val{search-runs} lần tìm kiếm của mỗi kích
thước, mỗi lần ghi một tệp nhật ký và một bản ghi JSON. Cuối cùng, \code{run_fleet_v3.bat reselect}
kết hợp
chúng thành các tập được bàn giao (Mục~\ref*{P-sec:fleet}). Trong kho mã HAN, \code{run_all.bat} chạy
thực nghiệm theo từng giai đoạn: kiểm tra dữ liệu, các đồ thị meta-path, hai lần rà soát, một
phép kiểm tra bộ nhớ ngắn, mô hình phổ biến, HAN, GAT, metapath2vec, các seed bổ sung và bảng kết
quả. Với tham số \code{dry}, \code{run_all.bat} in ra mọi lệnh mà không chạy. Mục~\ref*{P-sec:setup} nêu thời gian chạy.

\paragraph{Tính lặp lại.} Tìm kiếm dùng các kernel GPU tất định, nên seed quyết định một lần chạy:
hai lần chạy thử với seed~1 trùng khớp về $\alpha$, về nhóm ứng viên cuối và ở cả \val{pilot-sets}
tập được tính điểm. Việc huấn luyện HAN không lặp lại chính xác đến từng bit trên GPU, nhưng các
khác biệt nhỏ. Một lần chạy lặp lại (\code{a1} với bốn meta-path, seed~1) dừng ở cùng epoch với
cùng epoch tốt nhất (\val{repeat-han-best}), BPR xác thực (\val{repeat-han-val}) và Hit@10 thử
nghiệm (\val{repeat-han-hit}). Trong các nhóm người dùng, Hit@10 của nó khác nhau nhiều nhất
\val{repeat-han-groupdiff}.

\paragraph{Bản ghi.} Kho mã LMSPS lưu nhật ký và bản ghi của mọi lần tìm kiếm
(\code{hgb/logs/Fleet v3}) cùng lần chạy thử và các lần chọn lại. Kho mã HAN lưu nhật ký của mọi
lần chạy mô hình và bảng kết quả (\code{output/AMAZON/v3}), nhật ký của phép kiểm tra giai đoạn khởi
động (Mục~\ref*{P-sec:res-sensitivity}) và các ghi chú rà soát cùng script của chúng (\code{review}),
trong đó có đầu ra của phép kiểm tra theo nhóm sản phẩm (Mục~\ref*{P-sec:res-weights}). Script
\code{paper/make_numbers.py} trong kho mã LMSPS tính các con số, bảng và hình của bài báo và của tài
liệu này từ các bản ghi đó, với tùy chọn \code{--lang vi} cho định dạng số của bản tiếng Việt.
Script này cũng kiểm tra rằng mọi lần chạy đều dùng cùng dữ liệu và hai lần chạy thử cho kết quả giống nhau.

\paragraph{Công cụ gợi ý.} Kho mã HAN có các công cụ để xem một mô hình đã huấn luyện gợi ý những
gì. Script \code{recommend.py} liệt kê các sản phẩm hàng đầu cho một người dùng và các sản phẩm
giống nhất với một sản phẩm. Các script \code{recommend_6path.bat} và \code{recommend_4path.bat}
chạy nó trên GAT và, với HAN và metapath2vec, trên một kết quả đại diện cho mỗi loại tập meta-path.
Script \code{recommend_api.py} phục vụ cùng các truy vấn đó qua HTTP.

\paragraph{Không có trong các kho mã.} Các kho mã không chứa: (1) tệp thô, vốn vẫn nằm trên
trang của SNAP, (2) các tệp trung gian của bước trích xuất và tệp nút của đồ thị, mà
\code{generate_data.bat} dựng lại, (3) các mô hình đã huấn luyện (tệp \code{.pt}), mà phép kiểm tra theo nhóm sản phẩm (Mục~\ref*{P-sec:res-weights}) và các
công cụ gợi ý cần đến.

% =============================================================================================
\section{Các hạn chế khác}
\label{app:limits}

Mục~\ref*{P-sec:conclusion} nêu các hạn chế chính và tóm tắt hai hạn chế nhỏ hơn. Chi tiết của hai
hạn chế này là:
\begin{itemize}
  \item \emph{Chi tiết huấn luyện.} BPR xác thực vẫn còn cải thiện với trung vị \val{han-gain-med}
        trong mười epoch cuối (nhiều nhất \val{han-gain-max}, riêng các lần chạy của tập tìm được
        nhiều nhất \val{searched-gain-max}), nên huấn luyện lâu hơn nhiều khả năng sẽ nới rộng hơn
        là thu hẹp các khoảng cách. Tìm kiếm đã chọn các tập của nó trên các cạnh xác thực mà HAN
        cũng dùng để dừng sớm, nên các tập chỉ được so sánh trên các cạnh thử nghiệm. Các đồ thị
        meta-path qua người dùng của HAN chứa chính các cạnh huấn luyện mà nó học để dự đoán. Tìm
        kiếm có hiệu chỉnh điều này, HAN thì không, vì phép hiệu chỉnh không giúp ích cho HAN trên
        một mẫu nhỏ.
  \item \emph{Chỉ liên kết với danh mục lá:} nếu liên kết mọi danh mục cấp trên, các ứng viên và
        lựa chọn của tìm kiếm sẽ khác (Mục~\ref{app:graph}).
\end{itemize}


% The reference list stays in English, as the works were published (N7).
\begin{otherlanguage*}{english}
\bibliographystyle{ACM-Reference-Format}
\bibliography{references,notes_extra,reading_list}
\end{otherlanguage*}

\end{document}
